% Generated mechanically from the frozen algebraic.tex target.
% Do not edit this generated file; edit the bound locale source instead.

% OK, start here.
%

\setcounter{chapter}{93}
\chapter{代数叠}



\phantomsection
\label{algebraic-section-phantom}


\section{引言}
\label{algebraic-section-introduction}

\noindent
我们在此定义代数叠，并给出一些非常初步的观察。一般原则是完全不预设
任何分离性条件，只在使引理、命题和定理成立或可证时加入必要条件。
因此，这里讨论的概念与文献中其他地方所用的概念略有不同，例如
\cite{LM-B}。

\medskip\noindent
本章并非代数叠的入门介绍。关于代数叠的非形式化讨论，请参阅
《代数叠导论》第
\ref{stacks-introduction-section-introduction} 节。


\section{约定}
\label{algebraic-section-conventions}

\noindent
本章采用的约定与《代数空间》一章相同。为方便起见，我们在此重述。

\medskip\noindent
我们在适当的大 fppf 位点 $\Sch_{fppf}$ 中工作，参见
《概形上的拓扑》定义 \ref{topologies-definition-big-fppf-site}。
因此，除非另有明确说明，所有概形都是 $\Sch_{fppf}$ 的对象。
第 \ref{algebraic-section-change-big-site} 节将讨论更换大 fppf 位点时
会发生哪些变化。

\medskip\noindent
我们始终相对于 $\Sch_{fppf}$ 中的一个基 $S$ 工作，并使用大 fppf
位点 $(\Sch/S)_{fppf}$，参见《概形上的拓扑》定义
\ref{topologies-definition-big-small-fppf}。取
$S = \Spec(\mathbf{Z})$ 即可恢复绝对情形。

\medskip\noindent
若 $U,T$ 是 $S$ 上的概形，则用 $U(T)$ 表示相应的
{\it 在} $S$ 上的 $T$-值点集。用公式写即
$U(T) = \Mor_S(T, U)$。

\medskip\noindent
注意，任意 fpqc 覆盖都是泛有效满态射，参见
《下降》引理 \ref{descent-lemma-fpqc-universal-effective-epimorphisms}。
因此，$\Sch_{fppf}$ 上的拓扑弱于典范拓扑，而所有可表预层都是层。








\section{记号}
\label{algebraic-section-notation}

\noindent
我们用字母 $S,T,U,V,X,Y$ 表示概形，用字母
$\mathcal{X}, \mathcal{Y}, \mathcal{Z}$ 表示
$(\Sch/S)_{fppf}$ 上的范畴（纤维化范畴、以群胚为纤维的
纤维化范畴、叠等等）。我们用小写字母 $f$、$g$ 表示
$(\Sch/S)_{fppf}$ 上的函子，例如
$f : \mathcal{X} \to \mathcal{Y}$。
我们用大写字母 $F$、$G$、$H$ 表示 $S$ 上的代数空间，更一般地，
表示 $(\Sch/S)_{fppf}$ 上的集合预层。（在后续各章中，
我们也将重新使用 $X$、$Y$ 等字母表示代数空间。）

\medskip\noindent
作出这些选择，是为了在奠定基础时清楚地区分本章中的各种对象。









\section{可表的群胚纤维化范畴}
\label{algebraic-section-representable}

\noindent
设 $S$ 是 $\Sch_{fppf}$ 中的概形。本章的基本研究对象是
群胚纤维化范畴
$p : \mathcal{X} \to (\Sch/S)_{fppf}$，参见《范畴》定义
\ref{categories-definition-fibred-groupoids}。我们常简称
“设 $\mathcal{X}$ 是 $(\Sch/S)_{fppf}$ 上的群胚纤维化范畴”
来表示这一情形。$(\Sch/S)_{fppf}$ 上群胚纤维化范畴之间的
$1$-态射 $\mathcal{X} \to \mathcal{Y}$，是
$(\Sch/S)_{fppf}$ 上群胚纤维化范畴所成 $2$-范畴中的
$1$-态射，参见《范畴》定义
\ref{categories-definition-categories-fibred-in-groupoids-over-C}。
它就是 $(\Sch/S)_{fppf}$ 上的函子
$\mathcal{X} \to \mathcal{Y}$。回忆，这实际上是一个
$(2, 1)$-范畴，并且所有 $2$-纤维积都存在。

\medskip\noindent
设 $\mathcal{X}$ 是 $(\Sch/S)_{fppf}$ 上的群胚纤维化范畴。
回忆，若存在概形 $U \in \Ob((\Sch/S)_{fppf})$ 以及范畴等价
$$
j : \mathcal{X} \longrightarrow (\Sch/U)_{fppf}
$$
（它们都是 $(\Sch/S)_{fppf}$ 上的范畴），则称 $\mathcal{X}$
{\it 可表}，参见《范畴》定义
\ref{categories-definition-representable-fibred-category}。
为区别于 $\mathcal{X}$ 可由代数空间表示的情形（见下文），
我们有时称 $\mathcal{X}$ {\it 可由概形表示}。

\medskip\noindent
若 $\mathcal{X},\mathcal{Y}$ 是群胚纤维化范畴，并分别由
$U,V$ 表示，则有
\begin{equation}
\label{algebraic-equation-morphisms-schemes}
\Mor_{\textit{Cat}/(\Sch/S)_{fppf}}(\mathcal{X}, \mathcal{Y})
\Big/
2\text{-同构}
=
\Mor_{\Sch/S}(U, V)
\end{equation}
参见《范畴》引理
\ref{categories-lemma-morphisms-representable-fibred-categories}。
更确切地，任意 $1$-态射 $\mathcal{X} \to \mathcal{Y}$
都给出态射 $U \to V$。反过来，给定 $S$ 上的概形态射
$U \to V$，存在一个给出 $U \to V$ 的 $1$-态射
$\phi : \mathcal{X} \to \mathcal{Y}$，且它在唯一
$2$-同构意义下唯一。






\section{2-Yoneda 引理}
\label{algebraic-section-2-yoneda}

\noindent
设 $U \in \Ob((\Sch/S)_{fppf})$，并设 $\mathcal{X}$ 是
$(\Sch/S)_{fppf}$ 上的群胚纤维化范畴。我们将频繁使用
$2$-Yoneda 引理，参见《范畴》引理
\ref{categories-lemma-yoneda-2category}。严格来说，它断言存在范畴等价
$$
\Mor_{\textit{Cat}/(\Sch/S)_{fppf}}(
(\Sch/U)_{fppf}, \mathcal{X})
\longrightarrow
\mathcal{X}_U, \quad
f \longmapsto f(U/U).
$$
也就是说，$1$-态射 $(\Sch/U)_{fppf} \to \mathcal{X}$
对应于纤维范畴 $\mathcal{X}_U$ 的对象 $x$。具体地，给定
$1$-态射 $f : (\Sch/U)_{fppf} \to \mathcal{X}$，可得到对象
$x = f(U/U) \in \Ob(\mathcal{X}_U)$。反过来，按《范畴》定义
\ref{categories-definition-pullback-functor-fibred-category}
为 $\mathcal{X}$ 选定拉回，再给定 $\mathcal{X}_U$ 的对象 $x$，
便得到函子 $(\Sch/U)_{fppf} \to \mathcal{X}$；它在对象上的规则为
$$
(\varphi : V \to U) \longmapsto \varphi^*x
$$
。沿用记号，我们仍用
$x : (\Sch/U)_{fppf} \to \mathcal{X}$ 表示该函子。
它确实满足 $x(U/U) = x$；此外，对任意满足 $f(U/U) = x$
的其他函子 $f$，存在唯一的 $2$-同构 $x \to f$。换言之，
对象 $x$ 在唯一 $2$-同构意义下唯一确定函子 $x$。

\medskip\noindent
下文将直接使用这一事实，不再另行说明。





\section{群胚纤维化范畴的可表态射}
\label{algebraic-section-representable-morphism}


\noindent
设 $\mathcal{X}$、$\mathcal{Y}$ 是 $(\Sch/S)_{fppf}$ 上的
群胚纤维化范畴。设 $f : \mathcal{X} \to \mathcal{Y}$ 是
{\it 可表 $1$-态射}，参见《范畴》定义
\ref{categories-definition-representable-map-categories-fibred-in-groupoids}。
这意味着，对每个 $U \in \Ob((\Sch/S)_{fppf})$ 以及任意
$y \in \Ob(\mathcal{Y}_U)$，$2$-纤维积
$(\Sch/U)_{fppf} \times_{y, \mathcal{Y}} \mathcal{X}$
作为 $(\Sch/U)_{fppf}$ 上的群胚纤维化范畴是可表的。
选取表示对象 $f_y : V_y \to U$ 以及等价
$$
(\Sch/V_y)_{fppf}
\longrightarrow
(\Sch/U)_{fppf} \times_{y, \mathcal{Y}} \mathcal{X}.
$$
态射 $f_y$ 对应于投影
$(\Sch/V_y)_{fppf} \to
(\Sch/U)_{fppf} \times_\mathcal{Y} \mathcal{Y}
\to (\Sch/U)_{fppf}$
；参见第 \ref{algebraic-section-representable} 节公式
(\ref{algebraic-equation-morphisms-schemes})。我们用下图表示这一情形：
\begin{equation}
\label{algebraic-equation-representable}
\vcenter{
\xymatrix{
V_y \ar@{~>}[r] \ar[d]_{f_y} &
(\Sch/V_y)_{fppf} \ar[d] \ar[r] &
\mathcal{X} \ar[d]^f \\
U \ar@{~>}[r] &
(\Sch/U)_{fppf} \ar[r]^-y &
\mathcal{Y}
}
}
\end{equation}
其中波浪箭头表示 $2$-Yoneda 嵌入。下面给出关于这一概念的若干引理，
它们在极其一般的情形下成立（即对任意具有纤维积的基范畴上的
群胚纤维化范畴都成立）。

\begin{lemma}
\label{algebraic-lemma-morphism-schemes-gives-representable-transformation}
设 $f : X \to Y$ 是 $(\Sch/S)_{fppf}$ 中的态射。则 $f$ 诱导的
$1$-态射
$$
(\Sch/X)_{fppf} \longrightarrow (\Sch/Y)_{fppf}
$$
是可表 $1$-态射。
\end{lemma}

\begin{proof}
这是形式的，只用到范畴 $(\Sch/S)_{fppf}$ 具有纤维积这一事实。
\end{proof}

\begin{lemma}
\label{algebraic-lemma-representable-morphism-equivalent}
设 $S$ 是 $\Sch_{fppf}$ 的对象。考虑
$(\Sch/S)_{fppf}$ 上群胚纤维化范畴的 $1$-态射所成的
$2$-交换图
$$
\xymatrix{
\mathcal{X}' \ar[r] \ar[d]_{f'} & \mathcal{X} \ar[d]^f \\
\mathcal{Y}' \ar[r] & \mathcal{Y}
}
$$
。假设水平箭头都是等价。则 $f$ 可表当且仅当 $f'$ 可表。
\end{lemma}

\begin{proof}
略。
\end{proof}

\begin{lemma}
\label{algebraic-lemma-composition-representable-transformations}
设 $S$ 是 $\Sch_{fppf}$ 中的概形。设
$\mathcal{X},\mathcal{Y},\mathcal{Z}$ 是
$(\Sch/S)_{fppf}$ 上的群胚纤维化范畴，并设
$f : \mathcal{X} \to \mathcal{Y}$、
$g : \mathcal{Y} \to \mathcal{Z}$ 是可表 $1$-态射。则
$$
g \circ f : \mathcal{X} \longrightarrow \mathcal{Z}
$$
是可表 $1$-态射。
\end{lemma}

\begin{proof}
这完全是形式的，并在任意范畴中成立。
\end{proof}

\begin{lemma}
\label{algebraic-lemma-base-change-representable-transformations}
设 $S$ 是 $\Sch_{fppf}$ 中的概形。设
$\mathcal{X},\mathcal{Y},\mathcal{Z}$ 是
$(\Sch/S)_{fppf}$ 上的群胚纤维化范畴。设
$f : \mathcal{X} \to \mathcal{Y}$ 是可表 $1$-态射，并设
$g : \mathcal{Z} \to \mathcal{Y}$ 是任意 $1$-态射。
考虑纤维积图
$$
\xymatrix{
\mathcal{Z} \times_{g, \mathcal{Y}, f} \mathcal{X} \ar[r]_-{g'} \ar[d]_{f'} &
\mathcal{X} \ar[d]^f \\
\mathcal{Z} \ar[r]^g & \mathcal{Y}
}
$$
则基变换 $f'$ 是可表 $1$-态射。
\end{lemma}

\begin{proof}
这完全是形式的，并在任意范畴中成立。
\end{proof}

\begin{lemma}
\label{algebraic-lemma-product-representable-transformations}
设 $S$ 是 $\Sch_{fppf}$ 中的概形。设
$\mathcal{X}_i,\mathcal{Y}_i$ 是 $(\Sch/S)_{fppf}$ 上的
群胚纤维化范畴，$i = 1, 2$。设
$f_i : \mathcal{X}_i \to \mathcal{Y}_i$（$i = 1, 2$）
是可表 $1$-态射。则
$$
f_1 \times f_2 :
\mathcal{X}_1 \times \mathcal{X}_2
\longrightarrow
\mathcal{Y}_1 \times \mathcal{Y}_2
$$
是可表 $1$-态射。
\end{lemma}

\begin{proof}
把 $f_1 \times f_2$ 写成复合
$\mathcal{X}_1 \times \mathcal{X}_2 \to
\mathcal{Y}_1 \times \mathcal{X}_2 \to
\mathcal{Y}_1 \times \mathcal{Y}_2$。
第一个箭头是 $f_1$ 沿映射
$\mathcal{Y}_1 \times \mathcal{X}_2 \to \mathcal{Y}_1$ 的基变换，
第二个箭头是 $f_2$ 沿映射
$\mathcal{Y}_1 \times \mathcal{Y}_2 \to \mathcal{Y}_2$ 的基变换。
因此，本引理是引理
\ref{algebraic-lemma-composition-representable-transformations} 与
\ref{algebraic-lemma-base-change-representable-transformations}
的形式推论。
\end{proof}



\section{分裂群胚纤维化范畴}
\label{algebraic-section-split}

\noindent
设 $S$ 是 $\Sch_{fppf}$ 中的概形。回忆，给定一个“群胚预层”
$$
F : (\Sch/S)_{fppf}^{opp} \longrightarrow \textit{Groupoids}
$$
，可得到 $(\Sch/S)_{fppf}$ 上的群胚纤维化范畴
$\mathcal{S}_F$，参见《范畴》例
\ref{categories-example-functor-groupoids}。与其中某个范畴同构（！）的
群胚纤维化范畴称为{\it 分裂群胚纤维化范畴}。
任意群胚纤维化范畴都等价于一个分裂的群胚纤维化范畴。

\medskip\noindent
若 $F$ 是集合预层，则 $\mathcal{S}_F$ 以集合为纤维，参见
《范畴》定义 \ref{categories-definition-category-fibred-sets}
与《范畴》例 \ref{categories-example-presheaf}。
规则 $F \mapsto \mathcal{S}_F$ 在某种意义下于预层上全忠实，
参见《范畴》引理 \ref{categories-lemma-2-category-fibred-sets}。
若 $F,G$ 是预层，则
$$
\mathcal{S}_{F \times G}
=
\mathcal{S}_F \times_{(\Sch/S)_{fppf}} \mathcal{S}_G
$$
；若 $F \to H$ 与 $G \to H$ 是集合预层的映射，则
$$
\mathcal{S}_{F \times_H G} =
\mathcal{S}_F \times_{\mathcal{S}_H} \mathcal{S}_G
$$
，其中右端都是 $2$-纤维积。这由定义立即可得，因为
$\mathcal{S}_F, \mathcal{S}_G, \mathcal{S}_H$ 的各纤维范畴
只有恒等态射。

\medskip\noindent
更特殊的情形是 $F = h_X$ 为可表预层。此时
$\mathcal{S}_{h_X} = (\Sch/X)_{fppf}$，参见《范畴》例
\ref{categories-example-fibred-category-from-functor-of-points}。

\medskip\noindent
下文将直接使用记号 $\mathcal{S}_F$，不再另行说明。




\section{可由代数空间表示的群胚纤维化范畴}
\label{algebraic-section-representable-by-algebraic-spaces}

\noindent
比可表性稍弱的概念是可由代数空间表示，本节对此加以讨论。
若我们使用代数空间的某个范畴 $\textit{Spaces}_{fppf}$
而非范畴 $\Sch_{fppf}$，或许可以避免这番讨论。然而，在我们看来，
把概形范畴视为定义代数叠各纤维范畴所用的“测试对象”的自然集合，
是很自然的。

\medskip\noindent
仿照《范畴》定义
\ref{categories-definition-representable-fibred-category}，
作如下定义。

\begin{definition}
\label{algebraic-definition-representable-by-algebraic-space}
设 $S$ 是 $\Sch_{fppf}$ 中的概形。若存在 $S$ 上的代数空间
$F$ 以及 $(\Sch/S)_{fppf}$ 上范畴的等价
$j : \mathcal{X} \to \mathcal{S}_F$，则称群胚纤维化范畴
$p : \mathcal{X} \to (\Sch/S)_{fppf}$
{\it 可由 $S$ 上的代数空间表示}。
\end{definition}

\noindent
遇到这种情形时，我们继续沿用记号而省略等价 $j$。
由上文形式地可知，若 $\mathcal{X}$（由概形）可表，则它可由
代数空间表示。下面是《范畴》引理
\ref{categories-lemma-characterize-representable-fibred-category}
的对应结论。

\begin{lemma}
\label{algebraic-lemma-characterize-representable-by-space}
设 $S$ 是 $\Sch_{fppf}$ 中的概形，并设
$p : \mathcal{X} \to (\Sch/S)_{fppf}$ 是群胚纤维化范畴。
则 $\mathcal{X}$ 可由 $S$ 上的代数空间表示，当且仅当满足以下条件：
\begin{enumerate}
\item $\mathcal{X}$ 以集合胚为纤维\footnote{这意味着它以群胚为纤维，
且各纤维范畴中的对象没有非平凡自同构；参见《范畴》定义
\ref{categories-definition-category-fibred-setoids}。}；
\item 预层 $U \mapsto \Ob(\mathcal{X}_U)/\!\!\cong$ 是代数空间。
\end{enumerate}
\end{lemma}

\begin{proof}
略；但可参见《范畴》引理
\ref{categories-lemma-characterize-representable-fibred-category}。
\end{proof}

\noindent
若 $\mathcal{X}, \mathcal{Y}$ 以群胚为纤维，并分别可由
$S$ 上的代数空间 $F,G$ 表示，则有
\begin{equation}
\label{algebraic-equation-morphisms-spaces}
\Mor_{\textit{Cat}/(\Sch/S)_{fppf}}(\mathcal{X}, \mathcal{Y})
\Big/
2\text{-同构}
=
\Mor_{\Sch/S}(F, G)
\end{equation}
参见《范畴》引理
\ref{categories-lemma-2-category-fibred-setoids}。更确切地，任意
$1$-态射 $\mathcal{X} \to \mathcal{Y}$ 都给出态射 $F \to G$。
反过来，给定 $S$ 上的层态射 $F \to G$，存在一个给出
$F \to G$ 的 $1$-态射 $\phi : \mathcal{X} \to \mathcal{Y}$，
且它在唯一 $2$-同构意义下唯一。



\section{可由代数空间表示的态射}
\label{algebraic-section-morphisms-representable-by-algebraic-spaces}

\noindent
仿照《范畴》定义
\ref{categories-definition-representable-map-categories-fibred-in-groupoids}，
作如下定义。

\begin{definition}
\label{algebraic-definition-representable-by-algebraic-spaces}
设 $S$ 是 $\Sch_{fppf}$ 中的概形。若对任意
$U \in \Ob((\Sch/S)_{fppf})$ 及任意
$y : (\Sch/U)_{fppf} \to \mathcal{Y}$，群胚纤维化范畴
$$
(\Sch/U)_{fppf} \times_{y, \mathcal{Y}} \mathcal{X}
$$
作为 $(\Sch/U)_{fppf}$ 上的范畴可由 $U$ 上的代数空间表示，
则称 $(\Sch/S)_{fppf}$ 上群胚纤维化范畴的 $1$-态射
$f : \mathcal{X} \to \mathcal{Y}$ {\it 可由代数空间表示}。
\end{definition}

\noindent
选取表示 $(\Sch/U)_{fppf} \times_{y, \mathcal{Y}} \mathcal{X}$
的 $U$ 上代数空间 $F_y$。可把 $F_y$ 看成 $S$ 上的代数空间，
它配备一个典范的 $S$ 上态射 $f_y : F_y \to U$；参见
《代数空间》第 \ref{spaces-section-change-base-scheme} 节。
相应图表为
\begin{equation}
\label{algebraic-equation-representable-by-algebraic-spaces}
\vcenter{
\xymatrix{
F_y \ar[d]_{f_y} &
(\Sch/U)_{fppf} \times_{y, \mathcal{Y}} \mathcal{X}
\ar@{~>}[l] \ar[d]_{\text{pr}_0} \ar[r]_-{\text{pr}_1} &
\mathcal{X} \ar[d]^f \\
U &
(\Sch/U)_{fppf} \ar@{~>}[l] \ar[r]^-y &
\mathcal{Y}
}
}
\end{equation}
其中波浪箭头表示如下构造：把以集合胚为纤维的叠送到其对象同构类所成的
相应层。右侧方块是 $2$-交换的，并且是 $2$-纤维积方块。

\medskip\noindent
下面是《范畴》引理
\ref{categories-lemma-criterion-representable-map-stack-in-groupoids}
的对应结论。

\begin{lemma}
\label{algebraic-lemma-criterion-map-representable-spaces-fibred-in-groupoids}
设 $S$ 是 $\Sch_{fppf}$ 中的概形。设
$f : \mathcal{X} \to \mathcal{Y}$ 是 $(\Sch/S)_{fppf}$ 上
群胚纤维化范畴的 $1$-态射。则 $f$ 可由代数空间表示的充要条件如下：
\begin{enumerate}
\item 对每个概形 $U/S$，纤维范畴之间的函子
$f_U : \mathcal{X}_U \longrightarrow \mathcal{Y}_U$ 是忠实的；
\item 对每个 $U$ 以及每个 $y \in \Ob(\mathcal{Y}_U)$，预层
$$
(h : V \to U)
\longmapsto
\{(x, \phi) \mid x \in \Ob(\mathcal{X}_V), \phi : h^*y \to f(x)\}/\cong
$$
是 $U$ 上的代数空间。
\end{enumerate}
这里我们为 $\mathcal{Y}$ 选定了拉回。
\end{lemma}

\begin{proof}
把《范畴》引理 \ref{categories-lemma-identify-fibre-product}
中对 $2$-纤维积
$(\Sch/U)_{fppf} \times_{y, \mathcal{Y}} \mathcal{X}$
的纤维范畴的描述，与引理
\ref{algebraic-lemma-characterize-representable-by-space} 结合即可。
\end{proof}

\noindent
下面给出关于这一概念的若干引理；它们在非常一般的情形下成立。

\begin{lemma}
\label{algebraic-lemma-representable-by-spaces-morphism-equivalent}
设 $S$ 是 $\Sch_{fppf}$ 的对象。考虑
$(\Sch/S)_{fppf}$ 上群胚纤维化范畴的 $1$-态射所成的
$2$-交换图
$$
\xymatrix{
\mathcal{X}' \ar[r] \ar[d]_{f'} & \mathcal{X} \ar[d]^f \\
\mathcal{Y}' \ar[r] & \mathcal{Y}
}
$$
。假设水平箭头都是等价。则 $f$ 可由代数空间表示，当且仅当
$f'$ 可由代数空间表示。
\end{lemma}

\begin{proof}
略。
\end{proof}

\begin{lemma}
\label{algebraic-lemma-morphism-spaces-gives-representable-by-spaces}
设 $S$ 是 $\Sch_{fppf}$ 的对象。设
$f : \mathcal{X} \to \mathcal{Y}$ 是 $S$ 上群胚纤维化范畴的
$1$-态射。若 $\mathcal{X}$ 与 $\mathcal{Y}$ 都可由
$S$ 上的代数空间表示，则 $1$-态射 $f$ 可由代数空间表示。
\end{lemma}

\begin{proof}
略。这只用到 $S$ 上代数空间的范畴具有纤维积这一事实；
参见《代数空间》引理 \ref{spaces-lemma-fibre-product-spaces}。
\end{proof}

\begin{lemma}
\label{algebraic-lemma-map-presheaves-representable-by-algebraic-spaces}
设 $S$ 是 $\Sch_{fppf}$ 的对象。设 $a : F \to G$ 是
$(\Sch/S)_{fppf}$ 上集合预层的映射，并用
$a' : \mathcal{S}_F  \to \mathcal{S}_G$
表示相应的集合纤维化范畴映射。则 $a$ 可由代数空间表示
（参见《自举》定义
\ref{bootstrap-definition-morphism-representable-by-spaces}），
当且仅当 $a'$ 可由代数空间表示。
\end{lemma}

\begin{proof}
略。
\end{proof}

\begin{lemma}
\label{algebraic-lemma-map-fibred-setoids-representable-algebraic-spaces}
设 $S$ 是 $\Sch_{fppf}$ 的对象。设
$f : \mathcal{X} \to \mathcal{Y}$ 是 $(\Sch/S)_{fppf}$ 上
集合胚纤维化范畴的 $1$-态射。令 $F$、$G$ 分别为如下预层：
它把 $T$ 分别送到 $\mathcal{X}_T$、$\mathcal{Y}_T$
的对象同构类集合。令 $a : F \to G$ 为与 $f$ 对应的预层映射。
则 $a$ 可由代数空间表示（参见《自举》定义
\ref{bootstrap-definition-morphism-representable-by-spaces}），
当且仅当 $f$ 可由代数空间表示。
\end{lemma}

\begin{proof}
略。提示：结合引理
\ref{algebraic-lemma-representable-by-spaces-morphism-equivalent} 与
\ref{algebraic-lemma-map-presheaves-representable-by-algebraic-spaces}。
\end{proof}

\begin{lemma}
\label{algebraic-lemma-base-change-representable-by-spaces}
设 $S$ 是 $\Sch_{fppf}$ 中的概形。设
$\mathcal{X}, \mathcal{Y}, \mathcal{Z}$ 是
$(\Sch/S)_{fppf}$ 上的群胚纤维化范畴。设
$f : \mathcal{X} \to \mathcal{Y}$ 是可由代数空间表示的
$1$-态射，并设 $g : \mathcal{Z} \to \mathcal{Y}$
是任意 $1$-态射。考虑纤维积图
$$
\xymatrix{
\mathcal{Z} \times_{g, \mathcal{Y}, f} \mathcal{X} \ar[r]_-{g'} \ar[d]_{f'} &
\mathcal{X} \ar[d]^f \\
\mathcal{Z} \ar[r]^g & \mathcal{Y}
}
$$
则基变换 $f'$ 是可由代数空间表示的 $1$-态射。
\end{lemma}

\begin{proof}
这是形式的。
\end{proof}

\begin{lemma}
\label{algebraic-lemma-base-change-by-space-representable-by-space}
设 $S$ 是 $\Sch_{fppf}$ 中的概形。设
$\mathcal{X}, \mathcal{Y}, \mathcal{Z}$ 是
$(\Sch/S)_{fppf}$ 上的群胚纤维化范畴。设
$f : \mathcal{X} \to \mathcal{Y}$、
$g : \mathcal{Z} \to \mathcal{Y}$ 是 $1$-态射。假设
\begin{enumerate}
\item $f$ 可由代数空间表示；
\item $\mathcal{Z}$ 可由 $S$ 上的代数空间表示。
\end{enumerate}
则 $2$-纤维积
$\mathcal{Z} \times_{g, \mathcal{Y}, f} \mathcal{X}$
可由代数空间表示。
\end{lemma}

\begin{proof}
这是《自举》引理
\ref{bootstrap-lemma-representable-by-spaces-over-space}
的改述。首先注意，
$\mathcal{Z} \times_{g, \mathcal{Y}, f} \mathcal{X}$
在 $(\Sch/S)_{fppf}$ 上以集合胚为纤维。因此，它等价于
$\mathcal{S}_F$，其中 $F$ 是 $(\Sch/S)_{fppf}$ 上的某个预层；
参见《范畴》引理 \ref{categories-lemma-setoid-fibres}。
再设 $G$ 是表示 $\mathcal{Z}$ 的代数空间。$1$-态射
$\mathcal{Z} \times_{g, \mathcal{Y}, f} \mathcal{X} \to \mathcal{Z}$
由引理 \ref{algebraic-lemma-base-change-representable-by-spaces}
可由代数空间表示。又由《范畴》引理
\ref{categories-lemma-2-category-fibred-setoids}，态射
$\mathcal{Z} \times_{g, \mathcal{Y}, f} \mathcal{X} \to \mathcal{Z}$
对应于态射 $F \to G$。由引理
\ref{algebraic-lemma-map-fibred-setoids-representable-algebraic-spaces}，
$F \to G$ 可由代数空间表示。因此，《自举》引理
\ref{bootstrap-lemma-representable-by-spaces-over-space}
说明 $F$ 是代数空间，如所需。
\end{proof}

\noindent
设 $S$、$\mathcal{X}$、$\mathcal{Y}$、$\mathcal{Z}$、$f$、$g$
如引理 \ref{algebraic-lemma-base-change-by-space-representable-by-space} 中所述。
设 $F$ 与 $G$ 是 $S$ 上的代数空间，其中 $F$ 表示
$\mathcal{Z} \times_{g, \mathcal{Y}, f} \mathcal{X}$，而 $G$ 表示
$\mathcal{Z}$。$1$-态射
$f' : \mathcal{Z} \times_{g, \mathcal{Y}, f} \mathcal{X} \to \mathcal{Z}$
由 (\ref{algebraic-equation-morphisms-spaces}) 对应于代数空间的态射
$f' : F \to G$。因此有图表
\begin{equation}
\label{algebraic-equation-representable-by-algebraic-spaces-on-space}
\vcenter{
\xymatrix{
F \ar[d]_{f'} &
\mathcal{Z} \times_{g, \mathcal{Y}, f} \mathcal{X}
\ar@{~>}[l] \ar[d] \ar[r] &
\mathcal{X} \ar[d]^f \\
G &
\mathcal{Z} \ar@{~>}[l] \ar[r]^-g &
\mathcal{Y}
}
}
\end{equation}
其中波浪箭头表示如下构造：把以集合胚为纤维的叠送到其对象同构类所成的
相应层。

\begin{lemma}
\label{algebraic-lemma-composition-representable-by-spaces}
设 $S$ 是 $\Sch_{fppf}$ 中的概形。设
$\mathcal{X}, \mathcal{Y}, \mathcal{Z}$ 是
$(\Sch/S)_{fppf}$ 上的群胚纤维化范畴。若
$f : \mathcal{X} \to \mathcal{Y}$、
$g : \mathcal{Y} \to \mathcal{Z}$ 是可由代数空间表示的
$1$-态射，则
$$
g \circ f : \mathcal{X} \longrightarrow \mathcal{Z}
$$
是可由代数空间表示的 $1$-态射。
\end{lemma}

\begin{proof}
这由引理 \ref{algebraic-lemma-base-change-by-space-representable-by-space}
推出。略去细节。
\end{proof}

\begin{lemma}
\label{algebraic-lemma-product-representable-by-spaces}
设 $S$ 是 $\Sch_{fppf}$ 中的概形。设
$\mathcal{X}_i, \mathcal{Y}_i$ 是 $(\Sch/S)_{fppf}$ 上的
群胚纤维化范畴，$i = 1, 2$。设
$f_i : \mathcal{X}_i \to \mathcal{Y}_i$，$i = 1, 2$，
是可由代数空间表示的 $1$-态射。则
$$
f_1 \times f_2 :
\mathcal{X}_1 \times \mathcal{X}_2
\longrightarrow
\mathcal{Y}_1 \times \mathcal{Y}_2
$$
是可由代数空间表示的 $1$-态射。
\end{lemma}

\begin{proof}
把 $f_1 \times f_2$ 写成复合
$\mathcal{X}_1 \times \mathcal{X}_2 \to
\mathcal{Y}_1 \times \mathcal{X}_2 \to
\mathcal{Y}_1 \times \mathcal{Y}_2$。第一个箭头是 $f_1$ 沿映射
$\mathcal{Y}_1 \times \mathcal{X}_2 \to \mathcal{Y}_1$ 的基变换，
第二个箭头是 $f_2$ 沿映射
$\mathcal{Y}_1 \times \mathcal{Y}_2 \to \mathcal{Y}_2$ 的基变换。
因此，本引理是引理
\ref{algebraic-lemma-composition-representable-by-spaces} 与
\ref{algebraic-lemma-base-change-representable-by-spaces} 的形式推论。
\end{proof}

\begin{lemma}
\label{algebraic-lemma-get-a-stack}
\begin{reference}
Matthew Emerton 在 2016 年 6 月 15 日的一封电子邮件中的引理
\end{reference}
设 $S$ 是 $\Sch_{fppf}$ 中的概形。设
$\mathcal{X} \to \mathcal{Z}$ 与 $\mathcal{Y} \to \mathcal{Z}$
是 $(\Sch/S)_{fppf}$ 上群胚纤维化范畴的 $1$-态射。
若 $\mathcal{X} \to \mathcal{Z}$ 可由代数空间表示，且
$\mathcal{Y}$ 是群胚叠，则
$\mathcal{X} \times_\mathcal{Z} \mathcal{Y}$ 是群胚叠。
\end{lemma}

\begin{proof}
态射可由代数空间表示这一性质在基变换下保持
（引理 \ref{algebraic-lemma-base-change-by-space-representable-by-space}）。
因此，对 $\mathcal{Y}$ 作基变换
$\mathcal{X} \times_\mathcal{Z} \mathcal{Y}$ 后，可归结到如下情形：
群胚纤维化范畴的态射 $\mathcal{X} \to \mathcal{Y}$
可由代数空间表示，且其靶是群胚叠；此时目标是证明
$\mathcal{X}$ 也是群胚叠。这由《叠》引理
\ref{stacks-lemma-relative-sheaf-over-stack-is-stack} 得出；
引理 \ref{algebraic-lemma-criterion-map-representable-spaces-fibred-in-groupoids}
保证了其假设成立。
\end{proof}







\section{可由代数空间表示的态射的性质}
\label{algebraic-section-representable-properties}

\noindent
下面的定义使这一说法得以成立。

\begin{definition}
\label{algebraic-definition-relative-representable-property}
设 $S$ 是 $\Sch_{fppf}$ 中的概形。设
$f : \mathcal{X} \to \mathcal{Y}$ 是 $(\Sch/S)_{fppf}$ 上
群胚纤维化范畴的 $1$-态射。假设 $f$ 可由代数空间表示。
设 $\mathcal{P}$ 是代数空间态射的一个性质，并且
\begin{enumerate}
\item 它在任意基变换下保持；
\item 它在基上是 fppf 局部的；参见《下降与代数空间》定义
\ref{spaces-descent-definition-property-morphisms-local}。
\end{enumerate}
此时，若对每个 $U \in \Ob((\Sch/S)_{fppf})$ 及任意
$y \in \mathcal{Y}_U$，所得的代数空间态射
$f_y : F_y \to U$（见图
(\ref{algebraic-equation-representable-by-algebraic-spaces})）都具有性质
$\mathcal{P}$，则称 $f$ {\it 具有性质 $\mathcal{P}$}。
\end{definition}

\noindent
需要特别注意，我们只把此定义用于在基变换下稳定、并且在靶的 fppf
拓扑中为局部的态射性质。这并非因为在其他情况下该定义没有意义，
而是因为针对所考虑的性质，我们可能希望给出另一种更合适的定义。

\begin{lemma}
\label{algebraic-lemma-property-morphism-equivalent}
设 $S$ 是 $\Sch_{fppf}$ 的对象。设 $\mathcal{P}$ 如定义
\ref{algebraic-definition-relative-representable-property} 中所述。考虑
$(\Sch/S)_{fppf}$ 上群胚纤维化范畴的 $1$-态射所成的
$2$-交换图
$$
\xymatrix{
\mathcal{X}' \ar[r] \ar[d]_{f'} & \mathcal{X} \ar[d]^f \\
\mathcal{Y}' \ar[r] & \mathcal{Y}
}
$$
。假设水平箭头都是等价，并且 $f$（等价地，$f'$）可由代数空间表示。
则 $f$ 具有 $\mathcal{P}$ 当且仅当 $f'$ 具有 $\mathcal{P}$。
\end{lemma}

\begin{proof}
由引理 \ref{algebraic-lemma-representable-by-spaces-morphism-equivalent}
可知该陈述有意义。略去证明。
\end{proof}

\noindent
下面作一个基本检验。

\begin{lemma}
\label{algebraic-lemma-map-presheaves-representable-by-spaces-transformation-property}
设 $S$ 是 $\Sch_{fppf}$ 中的概形。设 $a : F \to G$ 是
$(\Sch/S)_{fppf}$ 上的预层映射。设 $\mathcal{P}$ 如定义
\ref{algebraic-definition-relative-representable-property} 中所述，并假设
$a$ 可由代数空间表示。则 $a : F \to G$ 具有性质
$\mathcal{P}$（参见《自举》定义
\ref{bootstrap-definition-property-transformation}），当且仅当
群胚纤维化范畴的对应态射 $\mathcal{S}_F \to \mathcal{S}_G$
具有性质 $\mathcal{P}$。
\end{lemma}

\begin{proof}
由引理 \ref{algebraic-lemma-map-presheaves-representable-by-algebraic-spaces}
可知本引理有意义。略去证明。
\end{proof}

\begin{lemma}
\label{algebraic-lemma-map-fibred-setoids-property}
设 $S$ 是 $\Sch_{fppf}$ 的对象。设 $\mathcal{P}$ 如定义
\ref{algebraic-definition-relative-representable-property} 中所述。
设 $f : \mathcal{X} \to \mathcal{Y}$ 是 $(\Sch/S)_{fppf}$ 上
集合胚纤维化范畴的 $1$-态射。令 $F$、$G$ 分别为如下预层：
它把 $T$ 分别送到 $\mathcal{X}_T$、$\mathcal{Y}_T$
的对象同构类集合。令 $a : F \to G$ 为与 $f$ 对应的预层映射。
则 $a$ 具有 $\mathcal{P}$ 当且仅当 $f$ 具有 $\mathcal{P}$。
\end{lemma}

\begin{proof}
由引理 \ref{algebraic-lemma-map-fibred-setoids-representable-algebraic-spaces}
可知本引理有意义。把引理 \ref{algebraic-lemma-property-morphism-equivalent}
与引理
\ref{algebraic-lemma-map-presheaves-representable-by-spaces-transformation-property}
结合即可得到结论。
\end{proof}

\begin{lemma}
\label{algebraic-lemma-composition-representable-transformations-property}
设 $S$ 是 $\Sch_{fppf}$ 中的概形。设
$\mathcal{X}$、$\mathcal{Y}$、$\mathcal{Z}$ 是
$(\Sch/S)_{fppf}$ 上的群胚纤维化范畴。设 $\mathcal{P}$
是定义 \ref{algebraic-definition-relative-representable-property}
所述的性质，并且对复合稳定。设
$f : \mathcal{X} \to \mathcal{Y}$、
$g : \mathcal{Y} \to \mathcal{Z}$ 是可由代数空间表示的
$1$-态射。若 $f$ 与 $g$ 具有性质 $\mathcal{P}$，则
$g \circ f : \mathcal{X} \to \mathcal{Z}$ 也具有该性质。
\end{lemma}

\begin{proof}
由引理 \ref{algebraic-lemma-composition-representable-by-spaces}
可知本引理有意义。略去证明。
\end{proof}

\begin{lemma}
\label{algebraic-lemma-base-change-representable-transformations-property}
设 $S$ 是 $\Sch_{fppf}$ 中的概形。设
$\mathcal{X}, \mathcal{Y}, \mathcal{Z}$ 是
$(\Sch/S)_{fppf}$ 上的群胚纤维化范畴。设 $\mathcal{P}$
为定义 \ref{algebraic-definition-relative-representable-property}
所述的性质。设 $f : \mathcal{X} \to \mathcal{Y}$
是可由代数空间表示的 $1$-态射，$g : \mathcal{Z} \to \mathcal{Y}$
是任意 $1$-态射。考虑 $2$-纤维积图
$$
\xymatrix{
\mathcal{Z} \times_{g, \mathcal{Y}, f} \mathcal{X} \ar[r]_-{g'} \ar[d]_{f'} &
\mathcal{X} \ar[d]^f \\
\mathcal{Z} \ar[r]^g & \mathcal{Y}
}
$$
若 $f$ 具有 $\mathcal{P}$，则基变换 $f'$ 也具有 $\mathcal{P}$。
\end{lemma}

\begin{proof}
由引理 \ref{algebraic-lemma-base-change-representable-by-spaces}
可知本引理有意义。略去证明。
\end{proof}

\begin{lemma}
\label{algebraic-lemma-descent-representable-transformations-property}
设 $S$ 是 $\Sch_{fppf}$ 中的概形。设
$\mathcal{X}, \mathcal{Y}, \mathcal{Z}$ 是
$(\Sch/S)_{fppf}$ 上的群胚纤维化范畴。设 $\mathcal{P}$
为定义 \ref{algebraic-definition-relative-representable-property}
所述的性质。设 $f : \mathcal{X} \to \mathcal{Y}$
是可由代数空间表示的 $1$-态射，$g : \mathcal{Z} \to \mathcal{Y}$
是任意 $1$-态射。考虑纤维积图
$$
\xymatrix{
\mathcal{Z} \times_{g, \mathcal{Y}, f} \mathcal{X} \ar[r]_-{g'} \ar[d]_{f'} &
\mathcal{X} \ar[d]^f \\
\mathcal{Z} \ar[r]^g & \mathcal{Y}
}
$$
假设对每个概形 $U$ 以及 $\mathcal{Y}_U$ 的对象 $x$，都存在
fppf 覆盖 $\{U_i \to U\}$，使得 $x|_{U_i}$ 属于函子
$g : \mathcal{Z}_{U_i} \to \mathcal{Y}_{U_i}$ 的本质像。
此时，若 $f'$ 具有 $\mathcal{P}$，则 $f$ 具有 $\mathcal{P}$。
\end{lemma}

\begin{proof}
略去证明。提示：比较《代数空间》引理
\ref{spaces-lemma-descent-representable-transformations-property}
的证明。
\end{proof}

\begin{lemma}
\label{algebraic-lemma-product-representable-transformations-property}
设 $S$ 是 $\Sch_{fppf}$ 中的概形。设 $\mathcal{P}$
为定义 \ref{algebraic-definition-relative-representable-property}
所述且对复合稳定的性质。设 $\mathcal{X}_i, \mathcal{Y}_i$
是 $(\Sch/S)_{fppf}$ 上的群胚纤维化范畴，$i = 1, 2$。
设 $f_i : \mathcal{X}_i \to \mathcal{Y}_i$，$i = 1, 2$，
是可由代数空间表示的 $1$-态射。若 $f_1$ 与 $f_2$ 具有性质
$\mathcal{P}$，则
$
f_1 \times f_2 :
\mathcal{X}_1 \times \mathcal{X}_2
\to
\mathcal{Y}_1 \times \mathcal{Y}_2
$
也具有该性质。
\end{lemma}

\begin{proof}
由引理 \ref{algebraic-lemma-product-representable-by-spaces}
可知本引理有意义。略去证明。
\end{proof}

\begin{lemma}
\label{algebraic-lemma-representable-transformations-property-implication}
设 $S$ 是 $\Sch_{fppf}$ 中的概形。设 $\mathcal{X}$、$\mathcal{Y}$
是 $(\Sch/S)_{fppf}$ 上的群胚纤维化范畴。设
$f : \mathcal{X} \to \mathcal{Y}$ 是可由代数空间表示的
$1$-态射。设 $\mathcal{P}$、$\mathcal{P}'$ 是定义
\ref{algebraic-definition-relative-representable-property} 所述的性质。
假设对代数空间的任意态射 $a : F \to G$，都有
$\mathcal{P}(a) \Rightarrow \mathcal{P}'(a)$。若 $f$ 具有性质
$\mathcal{P}$，则 $f$ 具有性质 $\mathcal{P}'$。
\end{lemma}

\begin{proof}
形式成立。
\end{proof}

\begin{lemma}
\label{algebraic-lemma-open-fibred-category-is-full}
设 $S$ 是 $\Sch_{fppf}$ 中的概形。设
$j : \mathcal X \to \mathcal Y$ 是 $(\Sch/S)_{fppf}$ 上
群胚纤维化范畴的 $1$-态射。假设 $j$ 可由代数空间表示且为单态射
（参见定义 \ref{algebraic-definition-relative-representable-property}
以及《下降与代数空间》引理
\ref{spaces-descent-lemma-descending-property-monomorphism}）。
则 $j$ 在各纤维范畴上全忠实。
\end{lemma}

\begin{proof}
由引理 \ref{algebraic-lemma-criterion-map-representable-spaces-fibred-in-groupoids}
已知 $j$ 在各纤维范畴上忠实。考虑概形 $U$、
$\mathcal{X}_U$ 的两个对象 $u, v$，以及 $\mathcal{Y}_U$ 中的同构
$t : j(u) \to j(v)$。我们需要在 $\mathcal{X}_U$ 中构造
$u$ 与 $v$ 之间的同构。由 $2$-Yoneda 引理（见第
\ref{algebraic-section-2-yoneda} 节），把 $u$、$v$ 看成 $1$-态射
$u, v : (\Sch/U)_{fppf} \to \mathcal{X}$
，并考虑 $2$-纤维积
$$
(\Sch/U)_{fppf} \times_{j \circ v, \mathcal{Y}} \mathcal{X}.
$$
根据假设，它可由 $U$ 上的代数空间 $F_{j \circ v}$ 表示，
而态射 $F_{j \circ v} \to U$ 是单态射。但
$(1_U, v, 1_{j(v)})$ 给出从 $(\Sch/U)_{fppf}$ 到上述
$2$-纤维积的 $1$-态射，故 $F_{j \circ v} = U$
（这里用到：若代数空间的单态射 $V \to U$ 有截面，则 $V = U$）。
因此，投影到第一坐标的 $1$-态射
$$
(\Sch/U)_{fppf} \times_{j \circ v, \mathcal{Y}} \mathcal{X}
\to (\Sch/U)_{fppf}
$$
是纤维范畴的等价。由于 $(1_U, u, t)$ 与
$(1_U, v, 1_{j(v)})$ 给出
$((\Sch/U)_{fppf} \times_{j \circ v, \mathcal{Y}}
\mathcal{X})_U$ 中具有相同第一坐标的两个对象，故它们之间必有
$2$-纤维积中的一个 $2$-态射。按定义，这就是满足
$j(\tilde t) = t$ 的态射 $\tilde t : u \to v$。
\end{proof}

\noindent
下面刻画对角态射可由代数空间表示的群胚纤维化范畴。

\begin{lemma}
\label{algebraic-lemma-representable-diagonal}
设 $S$ 是 $\Sch_{fppf}$ 中的概形。设 $\mathcal{X}$ 是
$(\Sch/S)_{fppf}$ 上的群胚纤维化范畴。以下各条件等价：
\begin{enumerate}
\item 对角态射 $\mathcal{X} \to \mathcal{X} \times \mathcal{X}$
可由代数空间表示；
\item 对 $S$ 上的每个概形 $U$ 以及任意
$x, y \in \Ob(\mathcal{X}_U)$，同构层
$\mathit{Isom}(x, y)$ 是 $U$ 上的代数空间；
\item 对 $S$ 上的每个概形 $U$ 以及任意
$x \in \Ob(\mathcal{X}_U)$，相应的 $1$-态射
$x : (\Sch/U)_{fppf} \to \mathcal{X}$ 可由代数空间表示；
\item 对 $S$ 上的每对概形 $T_1, T_2$ 以及任意
$x_i \in \Ob(\mathcal{X}_{T_i})$，$i = 1, 2$，$2$-纤维积
$(\Sch/T_1)_{fppf} \times_{x_1, \mathcal{X}, x_2}
(\Sch/T_2)_{fppf}$
可由代数空间表示；
\item 对 $(\Sch/S)_{fppf}$ 上的每个可表群胚纤维化范畴
$\mathcal{U}$，每个 $1$-态射 $\mathcal{U} \to \mathcal{X}$
都可由代数空间表示；
\item 对 $(\Sch/S)_{fppf}$ 上的每对可表群胚纤维化范畴
$\mathcal{T}_1, \mathcal{T}_2$，以及任意 $1$-态射
$x_i : \mathcal{T}_i \to \mathcal{X}$，$i = 1, 2$，$2$-纤维积
$\mathcal{T}_1 \times_{x_1, \mathcal{X}, x_2} \mathcal{T}_2$
可由代数空间表示；
\item 对 $(\Sch/S)_{fppf}$ 上可由代数空间表示的每个
群胚纤维化范畴 $\mathcal{U}$，每个 $1$-态射
$\mathcal{U} \to \mathcal{X}$ 都可由代数空间表示；
\item 对 $(\Sch/S)_{fppf}$ 上可由代数空间表示的每对
群胚纤维化范畴 $\mathcal{T}_1, \mathcal{T}_2$，以及任意
$1$-态射 $x_i : \mathcal{T}_i \to \mathcal{X}$，$2$-纤维积
$\mathcal{T}_1 \times_{x_1, \mathcal{X}, x_2} \mathcal{T}_2$
可由代数空间表示。
\end{enumerate}
\end{lemma}

\begin{proof}
(1) 与 (2) 的等价性由《叠》引理
\ref{stacks-lemma-isom-as-2-fibre-product} 和定义得出。
下面证明 (1) 与 (3) 等价。记基范畴为
$\mathcal{C} = (\Sch/S)_{fppf}$。我们将使用类似的《范畴》引理
\ref{categories-lemma-representable-diagonal-groupoids}
的证明中的若干观察。我们用符号 $\cong$ 表示“在
$\mathcal{C} = (\Sch/S)_{fppf}$ 上群胚纤维化范畴的等价”。
假设 (1) 成立。给定 (3) 中的 $U$ 与 $x$。对任意概形 $V$ 以及
$y \in \Ob(\mathcal{X}_V)$，有（比较上述引用）
$$
\mathcal{C}/U
\times_{x, \mathcal{X}, y}
\mathcal{C}/V
\cong
(\mathcal{C}/U \times_S V)
\times_{(x, y), \mathcal{X} \times \mathcal{X}, \Delta}
\mathcal{X}
$$
，它根据假设可由代数空间表示。反过来，假设 (3) 成立。考虑
$S$ 上任意概形 $U$，以及 $\mathcal{X}$ 在 $U$ 上的一对对象
$(x, x')$。我们需要证明
$\mathcal{X} \times_{\Delta, \mathcal{X} \times \mathcal{X}, (x, x')} U$
可由代数空间表示。这很清楚，因为（比较上述引用）
$$
\mathcal{X}
\times_{\Delta, \mathcal{X} \times \mathcal{X}, (x, x')}
\mathcal{C}/U
\cong
(\mathcal{C}/U \times_{x, \mathcal{X}, x'} \mathcal{C}/U)
\times_{\mathcal{C}/U \times_S U, \Delta}
\mathcal{C}/U
$$
，而右端根据假设以及如下事实可由代数空间表示：$S$ 上代数空间的范畴
具有纤维积，并且包含 $U$ 与 $S$。

\medskip\noindent
等价关系 (3) $\Leftrightarrow$ (4)、(5) $\Leftrightarrow$ (6)
以及 (7) $\Leftrightarrow$ (8) 都是形式的。等价关系
(3) $\Leftrightarrow$ (5) 与 (4) $\Leftrightarrow$ (6)
由引理 \ref{algebraic-lemma-representable-by-spaces-morphism-equivalent} 得出。
假设 (3) 成立，并设 $\mathcal{U} \to \mathcal{X}$ 如 (7) 中所述。
为证明 (7)，必须说明：对每个概形 $V$ 以及每个 $1$-态射
$y : (\Sch/V)_{fppf} \to \mathcal{X}$，$2$-纤维积
$(\Sch/V)_{fppf} \times_{y, \mathcal{X}} \mathcal{U}$
可由代数空间表示。性质 (3) 告诉我们 $y$ 可由代数空间表示，
故引理 \ref{algebraic-lemma-base-change-by-space-representable-by-space}
给出所需结论。最后，(7) 直接推出 (3)。
\end{proof}

\noindent
在该引理的情形下，对其中任意 $1$-态射
$x : (\Sch/U)_{fppf} \to \mathcal{X}$，以及定义
\ref{algebraic-definition-relative-representable-property} 所述的任意性质，
说 $x$ 具有性质 $\mathcal{P}$ 都是有意义的。
特别地，这适用于 $\mathcal{P} = $“满射”、
$\mathcal{P} = $“光滑”以及 $\mathcal{P} = $“\'etale”；
参见《下降与代数空间》引理
\ref{spaces-descent-lemma-descending-property-surjective}、
\ref{spaces-descent-lemma-descending-property-smooth} 与
\ref{spaces-descent-lemma-descending-property-etale}。
在下文定义代数叠时，我们将使用这三种情形。








\section{群胚叠}
\label{algebraic-section-stacks}

\noindent
设 $S$ 是 $\Sch_{fppf}$ 中的概形。回忆，
$(\Sch/S)_{fppf}$ 上的范畴
$p : \mathcal{X} \to (\Sch/S)_{fppf}$
称为{\it 群胚叠}（参见《叠》定义
\ref{stacks-definition-stack-in-groupoids}），当且仅当
\begin{enumerate}
\item $p : \mathcal{X} \to (\Sch/S)_{fppf}$ 在
$(\Sch/S)_{fppf}$ 上以群胚为纤维；
\item 对所有 $U \in \Ob((\Sch/S)_{fppf})$ 以及所有
$x, y\in \Ob(\mathcal{X}_U)$，预层 $\mathit{Isom}(x, y)$
是位点 $(\Sch/U)_{fppf}$ 上的层；
\item 对 $(\Sch/S)_{fppf}$ 中的所有覆盖
$\mathcal{U} = \{U_i \to U\}$，关于 $\mathcal{U}$ 的所有
下降数据 $(x_i, \phi_{ij})$ 都是有效的。
\end{enumerate}
关于例子，参见《叠的例子》第
\ref{examples-stacks-section-examples-stacks-in-groupoids}
节及以下各节。










\section{代数叠}
\label{algebraic-section-algebraic-stacks}

\noindent
下面给出代数叠的定义。注意，条件 (2) 使得第 (3) 项中
$(\Sch/U)_{fppf} \to \mathcal{X}$ 光滑且满这一条件有意义；
参见引理 \ref{algebraic-lemma-representable-diagonal} 后的讨论。

\begin{definition}
\label{algebraic-definition-algebraic-stack}
设 $S$ 是 $\Sch_{fppf}$ 中的基概形。{\it $S$ 上的代数叠}
是范畴
$$
p : \mathcal{X} \to (\Sch/S)_{fppf}
$$
，它位于 $(\Sch/S)_{fppf}$ 上，并具有以下性质：
\begin{enumerate}
\item 范畴 $\mathcal{X}$ 是 $(\Sch/S)_{fppf}$ 上的群胚叠。
\item 对角态射
$\Delta : \mathcal{X} \to \mathcal{X} \times \mathcal{X}$
可由代数空间表示。
\item 存在概形 $U \in \Ob((\Sch/S)_{fppf})$ 以及满且光滑的
$1$-态射 $(\Sch/U)_{fppf} \to \mathcal{X}$
\footnote{在后续各章中，我们将遵循文献中的惯例，把它简记为
$U \to \mathcal{X}$。另一种同样合适的表述是：存在一个可表的
群胚纤维化范畴 $\mathcal{U}$ 以及满且光滑的 $1$-态射
$\mathcal{U} \to \mathcal{X}$。}。
\end{enumerate}
\end{definition}

\noindent
这一定义与文献中的其他定义有若干差别。

\medskip\noindent
第一，我们要求 $\mathcal{X}$ 在 fppf 拓扑中为群胚叠，而许多文献
使用 \'etale 拓扑。在我们看来，fppf 拓扑似乎是自然的工作拓扑。
最终所得的代数叠 $2$-范畴仍然相同。《可表性判据》第
\ref{criteria-section-stacks-etale} 节对此作了说明。

\medskip\noindent
第二，我们只要求 $\mathcal{X}$ 的对角态射可由代数空间表示，
而大多数文献还施加其他条件。我们的观点是：先尝试只在
$\mathcal{X}$ 的对角态射可由代数空间表示这一假设下证明随后的一些
结果，并仅在必要处加入额外假设。这样还有一个好处：任意代数空间
（定义见《代数空间》定义 \ref{spaces-definition-algebraic-space}）
都给出一个代数叠。

\medskip\noindent
第三，一些论文要求存在概形 $U$ 以及满且 \'etale 的态射
$U \to \mathcal{X}$。在首次引入代数叠的奠基性论文 \cite{DM} 中，
Deligne 与 Mumford 采用了这一定义，并证明亏格 $g > 1$ 的稳定曲线
模叠是具有概形 \'etale 覆盖的代数叠。Michael Artin（参见
\cite{ArtinVersal}）认识到，关于代数叠的许多自然结果都能推广到
只假设存在概形光滑覆盖的情形。因此我们作出上述选择。为区分这两种
情形，文献中使用“Deligne--Mumford 叠”和“Artin 叠”这两个术语。
我们把“Artin 叠”一词保留到以后使用（此处插入未来引用），并继续使用
“代数叠”；但用“Deligne--Mumford 叠”指具有概形 \'etale 覆盖的
代数叠。

\begin{definition}
\label{algebraic-definition-deligne-mumford}
设 $S$ 是 $\Sch_{fppf}$ 中的概形。设 $\mathcal{X}$ 是
$S$ 上的代数叠。若存在概形 $U$ 以及满 \'etale 态射
$(\Sch/U)_{fppf} \to \mathcal{X}$，则称 $\mathcal{X}$
为{\it Deligne--Mumford 叠}。
\end{definition}

\noindent
稍后我们将比较这里的 Deligne--Mumford 叠概念与 Deligne 和 Mumford
论文中定义的概念（见此处插入未来引用）。

\medskip\noindent
$S$ 上的代数叠所成范畴构成一个 $2$-范畴。下面给出精确定义。

\begin{definition}
\label{algebraic-definition-morphism-algebraic-stacks}
设 $S$ 是 $\Sch_{fppf}$ 中的概形。{\it $S$ 上代数叠的
$2$-范畴}是 $(\Sch/S)_{fppf}$ 上群胚纤维化范畴所成
$2$-范畴（参见《范畴》定义
\ref{categories-definition-categories-fibred-in-groupoids-over-C}）
的如下子 $2$-范畴：
\begin{enumerate}
\item 其对象是 $(\Sch/S)_{fppf}$ 上那些构成 $S$ 上代数叠的
群胚纤维化范畴。
\item 其 $1$-态射 $f : \mathcal{X} \to \mathcal{Y}$ 是
$(\Sch/S)_{fppf}$ 上范畴的任意函子，如《范畴》定义
\ref{categories-definition-categories-over-C}。
\item 其 $2$-态射是 $(\Sch/S)_{fppf}$ 上函子之间的变换，
如《范畴》定义 \ref{categories-definition-categories-over-C}。
\end{enumerate}
\end{definition}

\noindent
换言之，这个 $2$-范畴是 $\textit{Cat}/(\Sch/S)_{fppf}$
中以代数叠为对象的满子 $2$-范畴。注意，每个 $2$-态射自动为同构。
因此，这实际上是一个 $(2, 1)$-范畴，而不仅是 $2$-范畴。

\medskip\noindent
稍后将看到（此处插入未来引用），这个 $2$-范畴具有 $2$-纤维积。

\medskip\noindent
与上面的注记类似，$S$ 上代数叠的 $2$-范畴是
$(\Sch/S)_{fppf}$ 上群胚纤维化范畴所成 $2$-范畴的满子
$2$-范畴。事实证明，它在等价下封闭。精确陈述如下。

\begin{lemma}
\label{algebraic-lemma-equivalent}
设 $S$ 是 $\Sch_{fppf}$ 中的概形。设 $\mathcal{X}$、$\mathcal{Y}$
是 $(\Sch/S)_{fppf}$ 上的范畴。假设 $\mathcal{X}$、$\mathcal{Y}$
作为 $(\Sch/S)_{fppf}$ 上的范畴等价。则 $\mathcal{X}$ 是代数叠
当且仅当 $\mathcal{Y}$ 是代数叠。类似地，$\mathcal{X}$ 是
Deligne--Mumford 叠当且仅当 $\mathcal{Y}$ 是 Deligne--Mumford 叠。
\end{lemma}

\begin{proof}
假设 $\mathcal{X}$ 是代数叠（相应地，是 Deligne--Mumford 叠）。
由《叠》引理 \ref{stacks-lemma-stack-in-groupoids-equivalent}，
$\mathcal{Y}$ 是 $\Sch_{fppf}$ 上的群胚叠。选取
$\Sch_{fppf}$ 上的等价 $f : \mathcal{X} \to \mathcal{Y}$。
这给出 $2$-交换图
$$
\xymatrix{
\mathcal{X} \ar[r]_f \ar[d]_{\Delta_\mathcal{X}} &
\mathcal{Y} \ar[d]^{\Delta_\mathcal{Y}} \\
\mathcal{X} \times \mathcal{X} \ar[r]^{f \times f} &
\mathcal{Y} \times \mathcal{Y}
}
$$
，其中水平箭头都是等价。根据引理
\ref{algebraic-lemma-representable-by-spaces-morphism-equivalent}，
这说明 $\Delta_\mathcal{Y}$ 可由代数空间表示。最后，设 $U$ 是
$S$ 上的概形，并设 $x : (\Sch/U)_{fppf} \to \mathcal{X}$
是满且光滑（相应地，\'etale）的 $1$-态射。考虑图表
$$
\xymatrix{
(\Sch/U)_{fppf} \ar[r]_{\text{id}} \ar[d]_x &
(\Sch/U)_{fppf} \ar[d]^{f \circ x} \\
\mathcal{X} \ar[r]^f &
\mathcal{Y}
}
$$
并应用引理 \ref{algebraic-lemma-property-morphism-equivalent}，可知
$f \circ x$ 满且光滑（相应地，\'etale），如所需。
\end{proof}




\section{代数叠与代数空间}
\label{algebraic-section-stacks-spaces}

\noindent
本节讨论若干蕴含代数叠为代数空间的简单判据。主要结果是：这恰好发生在
各纤维范畴的对象没有非平凡自同构时。这并非显然！在证明之前，
先作一个基本检验。

\begin{lemma}
\label{algebraic-lemma-representable-algebraic}
设 $S$ 是 $\Sch_{fppf}$ 中的概形。
\begin{enumerate}
\item 若群胚纤维化范畴
$p : \mathcal{X} \to (\Sch/S)_{fppf}$
可由代数空间表示，则它是 Deligne--Mumford 叠。
\item 若 $F$ 是 $S$ 上的代数空间，则相应的群胚纤维化范畴
$p : \mathcal{S}_F \to (\Sch/S)_{fppf}$
是 Deligne--Mumford 叠。
\item 若 $X \in \Ob((\Sch/S)_{fppf})$，则
$(\Sch/X)_{fppf} \to (\Sch/S)_{fppf}$
是 Deligne--Mumford 叠。
\end{enumerate}
\end{lemma}

\begin{proof}
(2) 蕴含 (3) 是清楚的。由引理 \ref{algebraic-lemma-equivalent}，
(1) 与 (2) 等价。因此只需证明 (2)。首先，由于 $F$ 是层，
$\mathcal{S}_F$ 是集合叠（《叠》引理
\ref{stacks-lemma-stack-in-setoids-characterize}），因而更是群胚叠。
其次，对角态射
$\mathcal{S}_F \to \mathcal{S}_F  \times \mathcal{S}_F$
就是由 $F$ 的对角态射产生的态射
$\mathcal{S}_F \to \mathcal{S}_{F \times F}$。故由引理
\ref{algebraic-lemma-morphism-spaces-gives-representable-by-spaces}，
它可由代数空间表示。事实上，它甚至（由概形）可表，因为代数空间的
对角态射可表，但我们不需要这一点。设 $U$ 是概形，并设
$h_U \to F$ 是满 \'etale 态射。可把它看作代数空间的满
\'etale 态射。于是由引理
\ref{algebraic-lemma-map-presheaves-representable-by-spaces-transformation-property}，
相应的 $1$-态射 $(\Sch/U)_{fppf} \to \mathcal{S}_F$
满且 \'etale。
\end{proof}

\noindent
下面的结果说明，惯性平凡的 Deligne--Mumford 叠“就是”代数空间。
下方更强的命题
\ref{algebraic-proposition-algebraic-stack-no-automorphisms}
将取代本引理；该命题说明这一结论对一般代数叠也成立……

\begin{lemma}
\label{algebraic-lemma-algebraic-stack-no-automorphisms}
设 $S$ 是 $\Sch_{fppf}$ 中的概形。设 $\mathcal{X}$ 是
$S$ 上的代数叠。以下各条件等价：
\begin{enumerate}
\item $\mathcal{X}$ 是 Deligne--Mumford 叠且为集合胚叠；
\item $\mathcal{X}$ 是 Deligne--Mumford 叠，且典范 $1$-态射
$\mathcal{I}_\mathcal{X} \to \mathcal{X}$ 是等价；
\item $\mathcal{X}$ 可由代数空间表示。
\end{enumerate}
\end{lemma}

\begin{proof}
(1) 与 (2) 的等价性由《叠》引理
\ref{stacks-lemma-characterize-stack-in-setoids} 得出。
(3) $\Rightarrow$ (1) 由引理 \ref{algebraic-lemma-representable-algebraic}
得出。最后，假设 (1) 成立。由《叠》引理
\ref{stacks-lemma-stack-in-setoids-characterize}，存在
$(\Sch/S)_{fppf}$ 上的层 $F$ 以及等价
$j : \mathcal{X} \to \mathcal{S}_F$。由引理
\ref{algebraic-lemma-map-presheaves-representable-by-algebraic-spaces}，
$\Delta_\mathcal{X}$ 可由代数空间表示意味着
$\Delta_F : F \to F \times F$ 可由代数空间表示。
设 $U$ 是概形，并设 $x : (\Sch/U)_{fppf} \to \mathcal{X}$
是满 \'etale 态射。复合
$j \circ x : (\Sch/U)_{fppf} \to \mathcal{S}_F$
对应于层的态射 $h_U \to F$。由《自举》引理
\ref{bootstrap-lemma-representable-diagonal}，该态射可由代数空间
表示。于是由引理 \ref{algebraic-lemma-map-fibred-setoids-property}，
$h_U \to F$ 满且 \'etale。最后应用《自举》定理
\ref{bootstrap-theorem-bootstrap}，可知 $F$ 是代数空间。
\end{proof}

\begin{proposition}
\label{algebraic-proposition-algebraic-stack-no-automorphisms}
设 $S$ 是 $\Sch_{fppf}$ 中的概形。设 $\mathcal{X}$ 是
$S$ 上的代数叠。以下各条件等价：
\begin{enumerate}
\item $\mathcal{X}$ 是集合胚叠；
\item 典范 $1$-态射 $\mathcal{I}_\mathcal{X} \to \mathcal{X}$
是等价；
\item $\mathcal{X}$ 可由代数空间表示。
\end{enumerate}
\end{proposition}

\begin{proof}
(1) 与 (2) 的等价性由《叠》引理
\ref{stacks-lemma-characterize-stack-in-setoids} 得出。
(3) $\Rightarrow$ (1) 由引理
\ref{algebraic-lemma-algebraic-stack-no-automorphisms} 得出。最后，假设 (1)
成立。由《叠》引理 \ref{stacks-lemma-stack-in-setoids-characterize}，
存在等价 $j : \mathcal{X} \to \mathcal{S}_F$，其中 $F$ 是
$(\Sch/S)_{fppf}$ 上的层。由引理
\ref{algebraic-lemma-map-presheaves-representable-by-algebraic-spaces}，
$\Delta_\mathcal{X}$ 可由代数空间表示意味着
$\Delta_F : F \to F \times F$ 可由代数空间表示。
设 $U$ 是概形，并设 $x : (\Sch/U)_{fppf} \to \mathcal{X}$
是满光滑态射。复合
$j \circ x : (\Sch/U)_{fppf} \to \mathcal{S}_F$
对应于层的态射 $h_U \to F$。由《自举》引理
\ref{bootstrap-lemma-representable-diagonal}，该态射可由代数空间
表示。因此由引理 \ref{algebraic-lemma-map-fibred-setoids-property}，
$h_U \to F$ 满且光滑。特别地，它是满、平坦且局部有限表示的
（这由引理
\ref{algebraic-lemma-representable-transformations-property-implication}
以及代数空间的光滑态射平坦且局部有限表示这一事实得出；参见
《代数空间的态射》引理
\ref{spaces-morphisms-lemma-smooth-locally-finite-presentation} 与
\ref{spaces-morphisms-lemma-smooth-flat}）。最后应用《自举》定理
\ref{bootstrap-theorem-final-bootstrap}，可知 $F$ 是代数空间。
\end{proof}






\section{代数叠的 2-纤维积}
\label{algebraic-section-2-fibre-products}

\noindent
代数叠的 $2$-范畴具有积与 $2$-纤维积。第一个引理其实是引理
\ref{algebraic-lemma-2-fibre-product} 的特例，但其证明稍为简单。

\begin{lemma}
\label{algebraic-lemma-product-spaces}
设 $S$ 是 $\Sch_{fppf}$ 中的概形。设 $\mathcal{X}$、$\mathcal{Y}$
是 $S$ 上的代数叠。则
$\mathcal{X} \times_{(\Sch/S)_{fppf}} \mathcal{Y}$
是代数叠，并且是 $S$ 上代数叠的 $2$-范畴中的积。
\end{lemma}

\begin{proof}
$\mathcal{X} \times_{(\Sch/S)_{fppf}} \mathcal{Y}$ 在 $T$ 上的
对象就是一对 $(x, y)$，其中 $x$ 是 $\mathcal{X}_T$ 的对象，
$y$ 是 $\mathcal{Y}_T$ 的对象。因此由定义立即可知，
$\mathcal{X} \times_{(\Sch/S)_{fppf}} \mathcal{Y}$ 是群胚叠。
若 $(x, y)$ 与 $(x', y')$ 是
$\mathcal{X} \times_{(\Sch/S)_{fppf}} \mathcal{Y}$
在 $T$ 上的两个对象，则
$$
\mathit{Isom}((x, y), (x', y')) =
\mathit{Isom}(x, x') \times \mathit{Isom}(y, y').
$$
于是，由引理 \ref{algebraic-lemma-representable-diagonal} 中的等价条件，
以及代数空间的范畴具有积这一事实，可知
$\mathcal{X} \times_{(\Sch/S)_{fppf}} \mathcal{Y}$ 的对角态射
可由代数空间表示。最后，设
$U, V \in \Ob((\Sch/S)_{fppf})$，并设 $x,y$ 是满光滑态射
$x : (\Sch/U)_{fppf} \to \mathcal{X}$,
$y : (\Sch/V)_{fppf} \to \mathcal{Y}$.
注意，
$$
(\Sch/U \times_S V)_{fppf} =
(\Sch/U)_{fppf}
\times_{(\Sch/S)_{fppf}} (\Sch/V)_{fppf}.
$$
因此，
$\mathcal{X} \times_{(\Sch/S)_{fppf}} \mathcal{Y}$ 在
$(\Sch/U \times_S V)_{fppf}$ 上的对象
$(\text{pr}_U^*x, \text{pr}_V^*y)$ 定义一个 $1$-态射
$$
(\Sch/U \times_S V)_{fppf}
\longrightarrow
\mathcal{X} \times_{(\Sch/S)_{fppf}} \mathcal{Y}
$$
。它是 $x$ 与 $y$ 的基变换的复合，故满且光滑；参见引理
\ref{algebraic-lemma-base-change-representable-transformations-property} 与
\ref{algebraic-lemma-composition-representable-transformations-property}。
因此，$\mathcal{X} \times_{(\Sch/S)_{fppf}} \mathcal{Y}$
确实是代数叠。略去它确实为积的验证。
\end{proof}

\begin{lemma}
\label{algebraic-lemma-2-fibre-product-general}
设 $S$ 是 $\Sch_{fppf}$ 中的概形。设 $\mathcal{Z}$ 是
$(\Sch/S)_{fppf}$ 上的群胚叠，且其对角态射可由代数空间表示。
设 $\mathcal{X}$、$\mathcal{Y}$ 是 $S$ 上的代数叠。
设 $f : \mathcal{X} \to \mathcal{Z}$、
$g : \mathcal{Y} \to \mathcal{Z}$ 是群胚叠的 $1$-态射。
则 $2$-纤维积
$\mathcal{X} \times_{f, \mathcal{Z}, g} \mathcal{Y}$ 是代数叠。
\end{lemma}

\begin{proof}
需要验证定义 \ref{algebraic-definition-algebraic-stack} 的条件 (1)、(2) 与
(3)。第一个条件由《叠》引理
\ref{stacks-lemma-2-product-stacks-in-groupoids} 得出。

\medskip\noindent
需要验证的第二个条件是 $\mathit{Isom}$ 层可由代数空间表示。
为此，设 $T$ 是 $S$ 上的概形，而 $u,v$ 是
$(\mathcal{X} \times_{f, \mathcal{Z}, g} \mathcal{Y})_T$
的对象。根据我们对 $2$-纤维积的构造（其源头可追溯到《范畴》引理
\ref{categories-lemma-2-product-categories-over-C}），可写成
$u = (x, y, \alpha)$ 与 $v = (x', y', \alpha')$，其中
$\alpha : f(x) \to g(y)$，$\alpha'$ 也类似。显然，
$$
\xymatrix{
\mathit{Isom}(u, v) \ar[d] \ar[rr] & &
\mathit{Isom}(y, y') \ar[d]^{\phi \mapsto g(\phi) \circ \alpha} \\
\mathit{Isom}(x, x') \ar[rr]^-{\psi \mapsto \alpha' \circ f(\psi)} & &
\mathit{Isom}(f(x), g(y'))
}
$$
是 $(\Sch/T)_{fppf}$ 上层的笛卡儿图。根据假设，层
$\mathit{Isom}(y, y')$、$\mathit{Isom}(x, x')$、
$\mathit{Isom}(f(x), g(y'))$ 都是代数空间（见引理
\ref{algebraic-lemma-representable-diagonal}），故
$\mathit{Isom}(u, v)$ 是代数空间。

\medskip\noindent
设 $U, V \in \Ob((\Sch/S)_{fppf})$，并设 $x,y$ 是满光滑态射
$x : (\Sch/U)_{fppf} \to \mathcal{X}$,
$y : (\Sch/V)_{fppf} \to \mathcal{Y}$.
考虑态射
$$
(\Sch/U)_{fppf}
\times_{f \circ x, \mathcal{Z}, g \circ y}
(\Sch/V)_{fppf}
\longrightarrow
\mathcal{X} \times_{f, \mathcal{Z}, g} \mathcal{Y}.
$$
由于 $\mathcal{Z}$ 的对角态射可由代数空间表示，该箭头的源可由
代数空间 $F$ 表示；见引理 \ref{algebraic-lemma-representable-diagonal}。
此外，该态射是 $x$ 与 $y$ 的基变换的复合，故满且光滑；参见引理
\ref{algebraic-lemma-base-change-representable-transformations-property} 与
\ref{algebraic-lemma-composition-representable-transformations-property}。
选取概形 $W$ 以及满 \'etale 态射 $W \to F$，可知上述
$1$-态射与相应 $1$-态射
$$
(\Sch/W)_{fppf}
\longrightarrow
(\Sch/U)_{fppf}
\times_{f \circ x, \mathcal{Z}, g \circ y}
(\Sch/V)_{fppf}
$$
的复合满且光滑，从而证明最后一个条件。
\end{proof}

\begin{lemma}
\label{algebraic-lemma-2-fibre-product}
设 $S$ 是 $\Sch_{fppf}$ 中的概形。设
$\mathcal{X}, \mathcal{Y}, \mathcal{Z}$ 是 $S$ 上的代数叠。
设 $f : \mathcal{X} \to \mathcal{Z}$、
$g : \mathcal{Y} \to \mathcal{Z}$ 是代数叠的 $1$-态射。
则 $2$-纤维积
$\mathcal{X} \times_{f, \mathcal{Z}, g} \mathcal{Y}$ 是代数叠。
它也是 $(\Sch/S)_{fppf}$ 上代数叠的 $2$-范畴中的 $2$-纤维积。
\end{lemma}

\begin{proof}
$\mathcal{X} \times_{f, \mathcal{Z}, g} \mathcal{Y}$ 是代数叠，
由更强的引理 \ref{algebraic-lemma-2-fibre-product-general} 得出。
而 $\mathcal{X} \times_{f, \mathcal{Z}, g} \mathcal{Y}$
是 $S$ 上代数叠的 $2$-范畴中的 $2$-纤维积，则形式地得自如下事实：
$S$ 上代数叠的 $2$-范畴是 $(\Sch/S)_{fppf}$ 上群胚叠所成
$2$-范畴的满子 $2$-范畴。
\end{proof}







\section{代数叠：重新整理}
\label{algebraic-section-overhaul}

\noindent
关于代数叠的一些基本结果。

\begin{lemma}
\label{algebraic-lemma-lift-morphism-presentations}
设 $S$ 是 $\Sch_{fppf}$ 中的概形。设
$f : \mathcal{X} \to \mathcal{Y}$ 是 $S$ 上代数叠的 $1$-态射。
设 $V \in \Ob((\Sch/S)_{fppf})$，并设
$y : (\Sch/V)_{fppf} \to \mathcal{Y}$ 满且光滑。
则存在对象 $U \in \Ob((\Sch/S)_{fppf})$ 以及 $2$-交换图
$$
\xymatrix{
(\Sch/U)_{fppf} \ar[r]_a \ar[d]_x &
(\Sch/V)_{fppf} \ar[d]^y \\
\mathcal{X} \ar[r]^f & \mathcal{Y}
}
$$
，其中 $x$ 满且光滑。
\end{lemma}

\begin{proof}
首先选取 $W \in \Ob((\Sch/S)_{fppf})$ 以及满光滑 $1$-态射
$z : (\Sch/W)_{fppf} \to \mathcal{X}$。由于 $\mathcal{Y}$
是代数叠，可选取等价
$$
j :
\mathcal{S}_F
\longrightarrow
(\Sch/W)_{fppf}
\times_{f \circ z, \mathcal{Y}, y}
(\Sch/V)_{fppf}
$$
，其中 $F$ 是代数空间。由引理
\ref{algebraic-lemma-base-change-representable-transformations-property}，
态射 $\mathcal{S}_F \to (\Sch/W)_{fppf}$ 作为 $y$ 的基变换
满且光滑。故由引理
\ref{algebraic-lemma-composition-representable-transformations-property}，
$\mathcal{S}_F \to \mathcal{X}$ 满且光滑。选取对象
$U \in \Ob((\Sch/S)_{fppf})$ 以及满 \'etale 态射 $U \to F$。
再次应用引理
\ref{algebraic-lemma-composition-representable-transformations-property}，
即得所需性质。
\end{proof}

\noindent
本引理是命题 \ref{algebraic-proposition-algebraic-stack-no-automorphisms}
的推广。

\begin{lemma}
\label{algebraic-lemma-characterize-representable-by-algebraic-spaces}
设 $S$ 是 $\Sch_{fppf}$ 中的概形。设
$f : \mathcal{X} \to \mathcal{Y}$ 是 $S$ 上代数叠的 $1$-态射。
以下各条件等价：
\begin{enumerate}
\item 对 $U \in \Ob((\Sch/S)_{fppf})$，函子
$f : \mathcal{X}_U \to \mathcal{Y}_U$ 是忠实的；
\item 函子 $f$ 是忠实的；
\item $f$ 可由代数空间表示。
\end{enumerate}
\end{lemma}

\begin{proof}
由群胚纤维化范畴的 $1$-态射的一般性质，(1) 与 (2) 等价；
参见《范畴》引理 \ref{categories-lemma-equivalence-fibred-categories}。
由引理 \ref{algebraic-lemma-criterion-map-representable-spaces-fibred-in-groupoids}，
(3) 蕴含 (2)。最后，假设 (2) 成立。设 $U$ 是概形，并设
$y \in \Ob(\mathcal{Y}_U)$。需要证明
$$
\mathcal{W} = (\Sch/U)_{fppf} \times_{y, \mathcal{Y}} \mathcal{X}
$$
可由 $U$ 上的代数空间表示。由于 $(\Sch/U)_{fppf}$ 是代数叠，
引理 \ref{algebraic-lemma-2-fibre-product} 说明 $\mathcal{W}$ 是代数叠。
另一方面，把 $\mathcal{W}$ 的对象明确描述为三元组
$(V, x, \alpha : y(V) \to f(x))$，再利用 $f$ 的忠实性，
可知 $\mathcal{W}$ 的纤维范畴都是集合胚。因此，命题
\ref{algebraic-proposition-algebraic-stack-no-automorphisms}
保证 $\mathcal{W}$ 可由代数空间表示。
\end{proof}

\begin{lemma}
\label{algebraic-lemma-smooth-surjective-morphism-implies-algebraic}
设 $S$ 是 $\Sch_{fppf}$ 中的概形。设
$u : \mathcal{U} \to \mathcal{X}$ 是 $(\Sch/S)_{fppf}$ 上
群胚叠的 $1$-态射。若
\begin{enumerate}
\item $\mathcal{U}$ 可由代数空间表示；
\item $u$ 可由代数空间表示、满且光滑，
\end{enumerate}
则 $\mathcal X$ 是 $S$ 上的代数叠。
\end{lemma}

\begin{proof}
需要证明对角态射
$\Delta : \mathcal{X} \to \mathcal{X} \times \mathcal{X}$
可由代数空间表示；参见定义 \ref{algebraic-definition-algebraic-stack}。
给定 $S$ 上的两个概形 $T_1$、$T_2$，用
$\mathcal{T}_i = (\Sch/T_i)_{fppf}$ 表示相应的可表纤维范畴。
设给定 $1$-态射 $f_i : \mathcal{T}_i \to \mathcal{X}$。
根据引理 \ref{algebraic-lemma-representable-diagonal}，只需证明 $2$-纤维积
$\mathcal{T}_1 \times_\mathcal{X} \mathcal{T}_2$
可由代数空间表示。由《叠》引理
\ref{stacks-lemma-2-fibre-product-stacks-in-setoids-over-stack-in-groupoids}，
它无论如何都是集合胚叠。因此，
$\mathcal{T}_1 \times_\mathcal{X} \mathcal{T}_2$
对应于 $(\Sch/S)_{fppf}$ 上的某个层 $F$；参见《叠》引理
\ref{stacks-lemma-stack-in-setoids-characterize}。
设 $U$ 是表示 $\mathcal{U}$ 的代数空间。根据假设，
$$
\mathcal{T}_i' = \mathcal{U} \times_{u, \mathcal{X}, f_i} \mathcal{T}_i
$$
可由 $S$ 上的代数空间 $T'_i$ 表示。因此，
$\mathcal{T}_1' \times_\mathcal{U} \mathcal{T}_2'$
可由代数空间 $T'_1 \times_U T'_2$ 表示。考虑交换图
$$
\xymatrix{
&
\mathcal{T}_1 \times_{\mathcal X} \mathcal{T}_2 \ar[rr]\ar'[d][dd] & &
\mathcal{T}_1 \ar[dd] \\
\mathcal{T}_1' \times_\mathcal{U} \mathcal{T}_2' \ar[ur]\ar[rr]\ar[dd] & &
\mathcal{T}_1' \ar[ur]\ar[dd] \\
&
\mathcal{T}_2 \ar'[r][rr] & &
\mathcal X \\
\mathcal{T}_2' \ar[rr]\ar[ur] & &
\mathcal{U} \ar[ur] }
$$
在此图中，底面、右面、背面和正面方块都是 $2$-纤维积。于是形式论证
说明，态射
$\mathcal{T}_1' \times_\mathcal{U} \mathcal{T}_2' \to
\mathcal{T}_1 \times_{\mathcal X} \mathcal{T}_2$
是 $\mathcal{U} \to \mathcal{X}$ 的“基变换”；更确切地，图
$$
\xymatrix{
\mathcal{T}_1' \times_\mathcal{U} \mathcal{T}_2' \ar[d] \ar[r] &
\mathcal{U} \ar[d] \\
\mathcal{T}_1 \times_{\mathcal X} \mathcal{T}_2 \ar[r] &
\mathcal{X}
}
$$
是 $2$-纤维积方块。因此，$T'_1 \times_U T'_2 \to F$
可由代数空间表示、光滑且满；参见引理
\ref{algebraic-lemma-map-fibred-setoids-representable-algebraic-spaces}、
\ref{algebraic-lemma-base-change-representable-by-spaces}、
\ref{algebraic-lemma-map-fibred-setoids-property} 与
\ref{algebraic-lemma-base-change-representable-transformations-property}。
于是由《自举》定理 \ref{bootstrap-theorem-final-bootstrap}，
$F$ 是代数空间，证毕。
\end{proof}

\noindent
引理 \ref{algebraic-lemma-smooth-surjective-morphism-implies-algebraic}
的一个应用是：代数叠上的代数空间仍是代数叠。这对应于《自举》引理
\ref{bootstrap-lemma-representable-by-spaces-over-space}。
事实上，只需假设态射 $\mathcal{X} \to \mathcal{Y}$
是“代数的”；参见《可表性判据》引理
\ref{criteria-lemma-algebraic-morphism-to-algebraic}。

\begin{lemma}
\label{algebraic-lemma-representable-morphism-to-algebraic}
设 $S$ 是 $\Sch_{fppf}$ 中的概形。设
$\mathcal{X} \to \mathcal{Y}$ 是 $(\Sch/S)_{fppf}$ 上群胚叠的态射。
假设
\begin{enumerate}
\item $\mathcal{X} \to \mathcal{Y}$ 可由代数空间表示；
\item $\mathcal{Y}$ 是 $S$ 上的代数叠。
\end{enumerate}
则 $\mathcal{X}$ 是 $S$ 上的代数叠。
\end{lemma}

\begin{proof}
设 $\mathcal{V} \to \mathcal{Y}$ 是从可表群胚叠到
$\mathcal{Y}$ 的满光滑 $1$-态射。由定义
\ref{algebraic-definition-algebraic-stack}，这样的态射存在。于是 $2$-纤维积
$\mathcal{U} = \mathcal{V} \times_{\mathcal Y} \mathcal X$
由引理 \ref{algebraic-lemma-base-change-by-space-representable-by-space}
可由代数空间表示。$1$-态射 $\mathcal{U} \to \mathcal X$
可由代数空间表示、光滑且满；参见引理
\ref{algebraic-lemma-base-change-representable-by-spaces} 与
\ref{algebraic-lemma-base-change-representable-transformations-property}。
由引理 \ref{algebraic-lemma-smooth-surjective-morphism-implies-algebraic}，
$\mathcal{X}$ 是代数叠。
\end{proof}

\begin{lemma}
\label{algebraic-lemma-open-fibred-category-is-algebraic}
\begin{reference}
Matthew Emerton 在日期为 2016 年 6 月 15 日的一封电子邮件中指出，
可以去掉 $j$ 为单态射这一假设
\end{reference}
设 $S$ 是 $\Sch_{fppf}$ 中的概形。设
$j : \mathcal X \to \mathcal Y$ 是 $(\Sch/S)_{fppf}$ 上
群胚纤维化范畴的 $1$-态射。假设 $j$ 可由代数空间表示。
若 $\mathcal{Y}$ 是群胚叠（相应地，代数叠），则
$\mathcal{X}$ 也是。
\end{lemma}

\begin{proof}
由引理 \ref{algebraic-lemma-representable-morphism-to-algebraic}，
关于代数叠的陈述可由关于群胚叠的陈述推出。若 $j$ 可由代数空间表示，
则 $j$ 在各纤维范畴上忠实，并且对每个 $U$ 与每个
$y \in \Ob(\mathcal{Y}_U)$，预层
$$
(h : V \to U)
\longmapsto
% P11-E0075：冻结源在此使用了未定义的 f(x)，而本引理只定义了 j。
\{(x, \phi) \mid x \in \Ob(\mathcal{X}_V), \phi : h^*y \to f(x)\}/\cong
$$
是 $U$ 上的代数空间；见引理
\ref{algebraic-lemma-criterion-map-representable-spaces-fibred-in-groupoids}。
特别地，该预层是层，故结论由《叠》引理
\ref{stacks-lemma-relative-sheaf-over-stack-is-stack} 得出。
\end{proof}



\section{从代数叠到表现}
\label{algebraic-section-stack-to-presentation}

\noindent
给定 $S$ 上的代数叠，可以得到 $S$ 上代数空间中的一个群胚，
其相应的商叠就是该代数叠。

\medskip\noindent
回忆，若 $(U, R, s, t, c)$ 是 $S$ 上代数空间中的群胚，则
$[U/R]$ 表示与该数据相应的商叠；参见《代数空间中的群胚》定义
\ref{spaces-groupoids-definition-quotient-stack}。
一般来说，$[U/R]$ {\bf 不是}代数叠。特别地，下一个引理中出现的
叠 $[U/R]$ 一般不是代数的。

\begin{lemma}
\label{algebraic-lemma-map-space-into-stack}
设 $S$ 是 $\Sch_{fppf}$ 中的概形。设 $\mathcal{X}$ 是
$S$ 上的代数叠。设 $\mathcal{U}$ 是 $S$ 上可由代数空间表示的
代数叠。设 $f : \mathcal{U} \to \mathcal{X}$ 是 1-态射。则
\begin{enumerate}
\item $2$-纤维积
$\mathcal{R} = \mathcal{U} \times_{f, \mathcal{X}, f} \mathcal{U}$
可由代数空间表示；
\item 存在典范等价
$$
\mathcal{U} \times_{f, \mathcal{X}, f} \mathcal{U}
\times_{f, \mathcal{X}, f} \mathcal{U} =
\mathcal{R} \times_{\text{pr}_1, \mathcal{U}, \text{pr}_0} \mathcal{R},
$$
\item 投影 $\text{pr}_{02}$ 通过 (2) 诱导 $1$-态射
$$
\text{pr}_{02} :
\mathcal{R} \times_{\text{pr}_1, \mathcal{U}, \text{pr}_0} \mathcal{R}
\longrightarrow
\mathcal{R}
$$
\item 设 $U$、$R$ 是分别表示 $\mathcal{U}, \mathcal{R}$
的代数空间，并设 $t, s : R \to U$ 与
$c : R \times_{s, U, t} R \to R$ 是分别对应于上述 $1$-态射
$\text{pr}_0, \text{pr}_1 : \mathcal{R} \to \mathcal{U}$
以及
$\text{pr}_{02} :
\mathcal{R} \times_{\text{pr}_1, \mathcal{U}, \text{pr}_0} \mathcal{R} \to
\mathcal{R}$ 的态射，则五元组 $(U, R, s, t, c)$ 是
$S$ 上代数空间中的群胚；
\item 态射 $f$ 诱导 $(\Sch/S)_{fppf}$ 上群胚叠的典范
$1$-态射
$f_{can} : [U/R] \to \mathcal{X}$
；
\item $1$-态射 $f_{can} : [U/R] \to \mathcal{X}$ 全忠实。
\end{enumerate}
\end{lemma}

\begin{proof}
(1) 的证明。按定义，$\Delta_\mathcal{X}$ 可由代数空间表示，
故应用引理 \ref{algebraic-lemma-representable-diagonal} 可知
$\mathcal{U} \to \mathcal{X}$ 可由代数空间表示。因此结论由引理
\ref{algebraic-lemma-base-change-by-space-representable-by-space} 得出。

\medskip\noindent
设 $T$ 是 $S$ 上的概形。根据 $2$-纤维积的构造（见《范畴》引理
\ref{categories-lemma-2-product-categories-over-C}），纤维范畴
$\mathcal{R}_T$ 的对象是三元组 $(a, b, \alpha)$，其中
$a, b \in \Ob(\mathcal{U}_T)$，而
$\alpha : f(a) \to f(b)$ 是纤维范畴 $\mathcal{X}_T$ 中的态射。

\medskip\noindent
(2) 的证明。反复应用《范畴》引理
\ref{categories-lemma-associativity-2-fibre-product} 与
\ref{categories-lemma-2-fibre-product-erase-factor}
即可得到该等价。把
$\mathcal{U} \times_\mathcal{X} \mathcal{U} \times_\mathcal{X} \mathcal{U}$
与
$(\mathcal{U} \times_\mathcal{X} \mathcal{U})
\times_\mathcal{X} \mathcal{U}$.
等同。若 $T$ 是 $S$ 上的概形，则该等价在 $T$ 上的纤维范畴中，
把左端对象 $((a, b, \alpha), c, \beta)$ 送到右端对象
$((a, b, \alpha), (b, c, \beta))$。

\medskip\noindent
(3) 的证明。《范畴》引理
\ref{categories-lemma-triple-2-fibre-product-pr02}
的证明构造了 $1$-态射 $\text{pr}_{02}$。根据上面对纤维范畴对象
的描述，$((a, b, \alpha), (b, c, \beta))$ 被送到
$(a, c, \beta \circ \alpha)$。

\medskip\noindent
遗憾的是，这与我们的群胚约定{\it 不相容}：我们总令
$j = (t, s) : R \to U$，并把 $R$ 的 $T$-值点 $r$
“看成”态射 $r : s(r) \to t(r)$。不过，这不影响 (4) 的证明，
因为群胚的反范畴仍是群胚。但在 (5) 的证明中，正是这一点导致
下方显示公式中出现逆。

\medskip\noindent
(4) 的证明。回忆，层 $U$ 与层
$T \mapsto \Ob(\mathcal{U}_T)/\!\cong$ 同构，$R$ 也类似；
参见引理 \ref{algebraic-lemma-characterize-representable-by-space}。
由《范畴》引理
\ref{categories-lemma-category-fibred-setoids-presheaves-products}，
这一描述与 $2$-纤维积相容，故类似地有
$\mathcal{R} \times_{\text{pr}_1, \mathcal{U}, \text{pr}_0} \mathcal{R}$
与 $R \times_{s, U, t} R$ 之间的对应。态射
$t, s : R \to U$ 与 $c : R \times_{s, U, t} R \to R$
来自一般等式 (\ref{algebraic-equation-morphisms-spaces})。明确地说，
% P11-E0078：冻结源在此重复 \text{pr}_0；后续结构所用三个投影为 \text{pr}_0、\text{pr}_1、\text{pr}_{02}。
这些映射是由 $\text{pr}_0$、$\text{pr}_0$、$\text{pr}_{02}$
作用于纤维范畴的对象同构类而产生的函子变换。因此，为证明所得结构
是代数空间中的群胚，只需证明对 $S$ 上的每个概形 $T$，结构
$$
(\Ob(\mathcal{U}_T)/\!\cong,
\Ob(\mathcal{R}_T)/\!\cong,
\text{pr}_1, \text{pr}_0, \text{pr}_{02})
$$
是群胚；这由上面对 $\mathcal{R}_T$ 的对象的描述立即可见。

\medskip\noindent
(5) 的证明。我们最终将应用《代数空间中的群胚》引理
\ref{spaces-groupoids-lemma-quotient-stack-2-coequalizer}
得到函子 $[U/R] \to \mathcal{X}$。考虑 $1$-态射
$f : \mathcal{U} \to \mathcal{X}$。由 $\mathcal{R}$
作为 $2$-纤维积的定义，有 $2$-箭头
$\tau : f \circ \text{pr}_1 \to f \circ \text{pr}_0$。
具体地，在 $\mathcal{R}$ 于 $T$ 上的对象 $(a, b, \alpha)$ 上，
它是映射 $\alpha^{-1} : b \to a$。我们断言
$$
\tau \circ \text{id}_{\text{pr}_{02}} =
(\tau \star \text{id}_{\text{pr}_0})
\circ
(\tau \star \text{id}_{\text{pr}_1}).
$$
该恒等式表示：给定
$\mathcal{R} \times_{\text{pr}_1, \mathcal{U}, \text{pr}_0} \mathcal{R}$
在 $T$ 上的对象 $((a, b, \alpha), (b, c, \beta))$，则
$$
\xymatrix{
c \ar[r]^{\beta^{-1}} & b \ar[r]^{\alpha^{-1}} & a
}
$$
% P11-E0079：冻结源所写箭头的源、靶与前一复合相反；此处保留原式待 canon 核验。
的复合等于箭头 $(\beta \circ \alpha)^{-1} : a \to c$。这显然成立，
故断言得证。由此可见，对结构
$(\mathcal{U}, \mathcal{R}, \text{pr}_0, \text{pr}_1, \text{pr}_{02})$
以及 $1$-态射 $f$ 和 $2$-态射 $\tau$，《代数空间中的群胚》引理
\ref{spaces-groupoids-lemma-quotient-stack-2-coequalizer}
的全部假设都成立。不过，为应用该引理，我们需要证明它对结构
$(\mathcal{S}_U, \mathcal{S}_R, s, t, c)$ 及适当态射也成立。

\medskip\noindent
此时应当存在某种一般的抽象论证，在二者之间传递这些数据，
但它似乎相当冗长。我们改用如下技巧。选取典范等价
$j : \mathcal{U} \to \mathcal{S}_U$ 的一个拟逆
$j^{-1} : \mathcal{S}_U \to \mathcal{U}$；该等价来自
$U(T) = \Ob(\mathcal{U}_T)/\!\!\cong$。这仅意味着：
对每个概形 $T/S$ 以及每个对象 $a \in \mathcal{U}_T$，
我们从其同构类中选定一个特定元素，即 $j^{-1}(j(a))$。
因此，利用 $j^{-1}$ 可把 $\mathcal{S}_U$ 看成 $\mathcal{U}$
的子范畴。选定这个子范畴后，可以考虑 $\mathcal{R}_T$
中满足如下条件的对象 $(a, b, \alpha)$：$a,b$ 是
$(\mathcal{S}_U)_T$ 的对象，即
$j^{-1}(j(a)) = a$ 且 $j^{-1}(j(b)) = b$。显然，它们构成
$\mathcal{R}$ 的一个子范畴，并通过典范等价
$\mathcal{R} \to \mathcal{S}_R$ 同构地映到 $\mathcal{S}_R$。
此外，这显然与形成 $2$-纤维积
$\mathcal{R} \times_{\text{pr}_1, \mathcal{U}, \text{pr}_0} \mathcal{R}$.
相容。因此，只需把 $f$ 限制到 $\mathcal{S}_U$，并把 $\tau$
限制为函子 $\mathcal{S}_R \to \mathcal{X}$ 之间的变换。
于是《代数空间中的群胚》引理
\ref{spaces-groupoids-lemma-quotient-stack-2-coequalizer}
中的显示等式显然成立，因为在限制到子范畴
$\mathcal{S}_{R \times_{s, U, t} R}$ 之前，它甚至已经作为函子
$\mathcal{R} \times_{\text{pr}_1, \mathcal{U}, \text{pr}_0} \mathcal{R}
\to \mathcal{X}$ 的变换等式成立。

\medskip\noindent
这证明《代数空间中的群胚》引理
\ref{spaces-groupoids-lemma-quotient-stack-2-coequalizer}
适用，从而得到所需的叠态射
$f_{can} : [U/R] \to \mathcal{X}$。我们简要说明在这个特殊情形中
$f_{can}$ 如何定义。在 $\mathcal{S}_U$ 于 $T$ 上的对象 $a$ 上，
有 $f_{can}(a) = f(a)$；这里通过上面选定的嵌入，把
$\mathcal{S}_U \subset \mathcal{U}$。若 $a,b$ 是
$\mathcal{S}_U$ 在 $T$ 上的对象，则按定义，$[U/R]$ 中的态射
$\varphi : a \to b$ 是 $\mathcal{R}$ 在 $T$ 上形如
$\varphi = (b, a, \alpha)$ 的对象。（注意次序交换。）
《代数空间中的群胚》引理
\ref{spaces-groupoids-lemma-quotient-stack-2-coequalizer}
的证明中的规则为
\begin{equation}
\label{algebraic-equation-on-morphisms}
f_{can}(\varphi) = \Big(f(a) \xrightarrow{\alpha^{-1}} f(b)\Big).
\end{equation}
(6) 的证明。$[U/R]$ 与 $\mathcal{X}$ 都是叠。因此，给定概形
$T/S$ 以及 $[U/R]$ 在 $T$ 上的对象 $a,b$，可得到
fppf 层的变换
$$
\mathit{Isom}(a, b) \longrightarrow \mathit{Isom}(f_{can}(a), f_{can}(b))
$$
，它位于 $(\Sch/T)_{fppf}$ 上。需要证明这是同构。可以在 $T$ 上
fppf 局部地工作，故可假设 $a,b$ 来自态射 $a, b : T \to U$。
通过上述嵌入 $\mathcal{S}_U \subset \mathcal{U}$，还可把
$a,b$ 看成 $\mathcal{U}$ 在 $T$ 上的对象。我们已在
《代数空间中的群胚》引理
\ref{spaces-groupoids-lemma-quotient-stack-morphisms} 中看到，
左端的层由代数空间
$$
R \times_{(t, s), U \times_S U, (b, a)} T
$$
表示于 $T$ 上。另一方面，由《叠》引理
\ref{stacks-lemma-isom-as-2-fibre-product}，右端等于与如下集合胚叠
相应的层：
$$
\mathcal{X}
\times_{\mathcal{X} \times \mathcal{X}, (f \circ b, f \circ a)} T =
\mathcal{X}
\times_{\mathcal{X} \times \mathcal{X}, (f, f)}
(\mathcal{U} \times \mathcal{U})
\times_{\mathcal{U} \times \mathcal{U}, (b, a)} T =
\mathcal{R}
\times_{(\text{pr}_0, \text{pr}_1), \mathcal{U} \times \mathcal{U}, (b, a)} T
$$
，它由上面显示的纤维积表示。至此已经证明两个
$\mathit{Isom}$ 层同构。由公式 (\ref{algebraic-equation-on-morphisms})，
我们的 $1$-态射 $f_{can} : [U/R] \to \mathcal{X}$
在 $\mathit{Isom}$ 层上诱导这个同构。
\end{proof}

\noindent
可以使用上述高度抽象的引理来构造表现。

\begin{lemma}
\label{algebraic-lemma-stack-presentation}
设 $S$ 是 $\Sch_{fppf}$ 中的概形。设 $\mathcal{X}$ 是
$S$ 上的代数叠。设 $U$ 是 $S$ 上的代数空间，并设
$f : \mathcal{S}_U \to \mathcal{X}$ 是满光滑态射。
把引理 \ref{algebraic-lemma-map-space-into-stack} 应用于 $U$ 与 $f$，
得到代数空间中的群胚 $(U, R, s, t, c)$ 以及
$f_{can} : [U/R] \to \mathcal{X}$。则
\begin{enumerate}
\item 态射 $s$、$t$ 光滑；
\item $1$-态射 $f_{can} : [U/R] \to \mathcal{X}$ 是等价。
\end{enumerate}
\end{lemma}

\begin{proof}
由引理 \ref{algebraic-lemma-property-morphism-equivalent} 与
\ref{algebraic-lemma-map-presheaves-representable-by-spaces-transformation-property}，
态射 $s,t$ 光滑。由于 $1$-态射 $f$ 光滑且满，显然，对任意概形
$T$ 以及任意对象 $a \in \Ob(\mathcal{X}_T)$，存在光滑满态射
$T' \to T$，使得 $a|_T'$ 来自 $[U/R]_{T'}$ 的对象。
由于 $f_{can} : [U/R] \to \mathcal{X}$ 全忠实，而两端的对象下降
数据都有效，故 $[U/R] \to \mathcal{X}$ 本质满；参见《叠》引理
\ref{stacks-lemma-characterize-essentially-surjective-when-ff}。
\end{proof}

\begin{remark}
\label{algebraic-remark-flat-fp-presentation}
若对引理 \ref{algebraic-lemma-stack-presentation} 中的态射
$f : \mathcal{S}_U \to \mathcal{X}$ 只假设其满、平坦且局部有限表示，
则 $f_{can} : [U/R] \to \mathcal{X}$ 仍是等价。此时态射
$s$、$t$ 平坦且局部有限表示，但一般而言当然不光滑。
\end{remark}

\noindent
引理 \ref{algebraic-lemma-stack-presentation} 提示如下定义。

\begin{definition}
\label{algebraic-definition-smooth-groupoid}
设 $S$ 是概形。设 $B$ 是 $S$ 上的代数空间。设
$(U, R, s, t, c)$ 是 $B$ 上代数空间中的群胚。若
$s, t : R \to U$ 是代数空间的光滑态射，则称
$(U, R, s, t, c)$ 为{\it 光滑群胚}\footnote{这个术语可能稍显
混淆：它并不蕴含 $[U/R]$ 在任何对象上光滑。}。
\end{definition}

\begin{definition}
\label{algebraic-definition-presentation}
设 $\mathcal{X}$ 是 $S$ 上的代数叠。$\mathcal{X}$ 的一个
{\it 表现}由 $S$ 上代数空间中的光滑群胚
$(U, R, s, t, c)$ 以及等价 $f : [U/R] \to \mathcal{X}$ 给出。
\end{definition}

\noindent
上文已经看到，每个代数叠都有表现。下一步将证明，$S$ 上代数空间中的
每个光滑群胚都给出一个代数叠。


\section{与光滑群胚相应的代数叠}
\label{algebraic-section-smooth-groupoid-gives-algebraic-stack}

\noindent
本节从代数空间中的光滑群胚出发，证明相应的商叠是代数叠。

\begin{lemma}
\label{algebraic-lemma-diagonal-quotient-stack}
设 $S$ 是 $\Sch_{fppf}$ 中的概形。设 $(U, R, s, t, c)$ 是
$S$ 上代数空间中的群胚。则 $[U/R]$ 的对角态射可由代数空间表示。
\end{lemma}

\begin{proof}
只需证明 $\mathit{Isom}$ 层是代数空间；见引理
\ref{algebraic-lemma-representable-diagonal}。这由《自举》引理
\ref{bootstrap-lemma-quotient-stack-isom} 得出。
\end{proof}

\begin{lemma}
\label{algebraic-lemma-smooth-quotient-smooth-presentation}
设 $S$ 是 $\Sch_{fppf}$ 中的概形。设 $(U, R, s, t, c)$ 是
$S$ 上代数空间中的光滑群胚。则态射
$\mathcal{S}_U \to [U/R]$ 光滑且满。
\end{lemma}

\begin{proof}
设 $T$ 是概形，并设 $x : (\Sch/T)_{fppf} \to [U/R]$
是 $1$-态射。需要证明投影
$$
\mathcal{S}_U \times_{[U/R]} (\Sch/T)_{fppf}
\longrightarrow
(\Sch/T)_{fppf}
$$
满且光滑。由引理 \ref{algebraic-lemma-diagonal-quotient-stack} 与
\ref{algebraic-lemma-representable-diagonal} 已知，左端可由代数空间 $F$
表示。因此，需要证明相应的代数空间态射 $F \to T$ 满且光滑。
我们研究的是在靶的 fppf 拓扑中为局部的代数空间态射性质，
故可以在 $T$ 上 fppf 局部地验证。按构造，存在 $T$ 的 fppf 覆盖
$\{T_i \to T\}$，使得 $x|_{(\Sch/T_i)_{fppf}}$ 来自态射
$x_i : T_i \to U$。（注意，$F \times_T T_i$ 表示 $2$-纤维积
$\mathcal{S}_U \times_{[U/R]} (\Sch/T_i)_{fppf}$，故一切都与
沿 $T_i \to T$ 的基变换相容。）因此可假设 $x$ 来自
$x : T \to U$。此时有
$$
\mathcal{S}_U \times_{[U/R]} (\Sch/T)_{fppf}
=
(\mathcal{S}_U \times_{[U/R]} \mathcal{S}_U)
\times_{\mathcal{S}_U} (\Sch/T)_{fppf}
=
\mathcal{S}_R \times_{\mathcal{S}_U} (\Sch/T)_{fppf}
$$
第一个等式由《范畴》引理
\ref{categories-lemma-2-fibre-product-erase-factor} 得出，
第二个等式由《代数空间中的群胚》引理
\ref{spaces-groupoids-lemma-quotient-stack-2-cartesian} 得出。
显然，最后一个 $2$-纤维积由代数空间
$F = R \times_{s, U, x} T$ 表示，而投影
$R \times_{s, U, x} T \to T$ 是代数空间的光滑态射
$s : R \to U$ 的基变换，故光滑。它也满，因为 $s$ 有截面
（即群胚的恒等元 $e : U \to R$）。引理得证。
\end{proof}

\noindent
下面是本节的主要结果。

\begin{theorem}
\label{algebraic-theorem-smooth-groupoid-gives-algebraic-stack}
设 $S$ 是 $\Sch_{fppf}$ 中的概形。设 $(U, R, s, t, c)$ 是
$S$ 上代数空间中的光滑群胚。则商叠 $[U/R]$ 是 $S$ 上的代数叠。
\end{theorem}

\begin{proof}
验证定义 \ref{algebraic-definition-algebraic-stack} 的三个条件。
按构造，$[U/R]$ 是群胚叠，这就是第一个条件。

\medskip\noindent
第二个条件由更强的引理 \ref{algebraic-lemma-diagonal-quotient-stack} 得出。

\medskip\noindent
最后，需要证明存在 $S$ 上的概形 $W$ 以及满光滑 $1$-态射
% P11-E0080：冻结源在此使用未定义的 \mathcal{X}；本定理上下文的商叠为 [U/R]。
$(\Sch/W)_{fppf} \longrightarrow \mathcal{X}$.
先选取 $W \in \Ob((\Sch/S)_{fppf})$ 以及满 \'etale 态射
$W \to U$。注意，这给出集合纤维化范畴的满 \'etale 态射
$\mathcal{S}_W \to \mathcal{S}_U$；见引理
\ref{algebraic-lemma-map-presheaves-representable-by-spaces-transformation-property}。
当然，$\mathcal{S}_W \to \mathcal{S}_U$ 也满且光滑；
见引理 \ref{algebraic-lemma-representable-transformations-property-implication}。
因此，结合引理 \ref{algebraic-lemma-smooth-quotient-smooth-presentation} 与
\ref{algebraic-lemma-composition-representable-transformations-property}，
$\mathcal{S}_W \to \mathcal{S}_U \to [U/R]$ 满且光滑。
\end{proof}





\section{更换大位点}
\label{algebraic-section-change-big-site}

\noindent
本节简要讨论更换大位点时发生的情形。要点是，我们总能任意扩大大位点；
因此，可以假设想要考虑的任意概形集合都包含在我们研究代数空间所用的
大 fppf 位点中。我们建议读者跳过本节。

\medskip\noindent
叠的拉回定义见《叠》第 \ref{stacks-section-inverse-image} 节。

\begin{lemma}
\label{algebraic-lemma-change-big-site}
设给定大位点 $\Sch_{fppf}$ 与 $\Sch'_{fppf}$。假设
$\Sch_{fppf}$ 包含于 $\Sch'_{fppf}$；参见《概形上的拓扑》第
\ref{topologies-section-change-alpha} 节。设 $S$ 是
$\Sch_{fppf}$ 的对象。令
$f : (\Sch'/S)_{fppf} \to (\Sch/S)_{fppf}$ 为与包含函子
$u : (\Sch/S)_{fppf} \to (\Sch'/S)_{fppf}$ 相应的位点态射。
设 $\mathcal{X}$ 是 $(\Sch/S)_{fppf}$ 上的群胚叠。
\begin{enumerate}
\item 若 $\mathcal{X}$ 可由某个
$X \in \Ob((\Sch/S)_{fppf})$ 表示，则
$f^{-1}\mathcal{X}$ 也可表；事实上，它由同一个概形 $X$ 表示，
只是现在将该概形看成 $(\Sch'/S)_{fppf}$ 的对象；
\item 若 $\mathcal{X}$ 可由代数空间
$F \in \Sh((\Sch/S)_{fppf})$ 表示，则
$f^{-1}\mathcal{X}$ 可由代数空间 $f^{-1}F$ 表示；
\item 若 $\mathcal{X}$ 是代数叠，则 $f^{-1}\mathcal{X}$ 是代数叠；
\item 若 $\mathcal{X}$ 是 Deligne--Mumford 叠，则
$f^{-1}\mathcal{X}$ 也是 Deligne--Mumford 叠。
\end{enumerate}
\end{lemma}

\begin{proof}
证明 (3)。由引理 \ref{algebraic-lemma-stack-presentation}，可对某个
代数空间中的光滑群胚 $(U, R, s, t, c)$ 写成
$\mathcal{X} = [U/R]$。由《代数空间中的群胚》引理
\ref{spaces-groupoids-lemma-quotient-stack-change-big-site}，
有 $f^{-1}[U/R] = [f^{-1}U/f^{-1}R]$。当然，
$(f^{-1}U, f^{-1}R, f^{-1}s, f^{-1}t, f^{-1}c)$
也是代数空间中的光滑群胚。因此 (3) 得证。

\medskip\noindent
其他情形 (1)、(2)、(4) 各自都意味着 $\mathcal{X}$ 具有某种特定类型
的表现 $[U/R]$，因而转化为
$f^{-1}\mathcal{X} = [f^{-1}U/f^{-1}R]$ 的同类型表现。
引理得证。
\end{proof}

\noindent
一般而言，大位点上的代数空间限制到小位点后未必仍是小位点上的
代数空间（仅基数方面的原因就足以导致失败）。因此，这类简单引理
只能用来扩大基范畴，绝不能用来缩小它。

\begin{lemma}
\label{algebraic-lemma-fully-faithful}
假设 $\Sch_{fppf}$ 包含于 $\Sch'_{fppf}$。设 $S$ 是
$\Sch_{fppf}$ 的对象。用 $\textit{Algebraic-Stacks}/S$
表示用 $\Sch_{fppf}$ 定义的 $S$ 上代数叠的 $2$-范畴。
类似地，用 $\textit{Algebraic-Stacks}'/S$ 表示用
$\Sch'_{fppf}$ 定义的 $S$ 上代数叠的 $2$-范畴。引理
\ref{algebraic-lemma-change-big-site} 的规则
$\mathcal{X} \mapsto f^{-1}\mathcal{X}$
定义 $2$-范畴的函子
$$
\textit{Algebraic-Stacks}/S \longrightarrow \textit{Algebraic-Stacks}'/S
$$
，并对 $\textit{Algebraic-Stacks}/S$ 的每对对象
$\mathcal{X}, \mathcal{Y}$ 定义态射范畴的等价
$$
\Mor_{\textit{Algebraic-Stacks}/S}(\mathcal{X}, \mathcal{Y})
\longrightarrow
\Mor_{\textit{Algebraic-Stacks}'/S}(f^{-1}\mathcal{X}, f^{-1}\mathcal{Y})
$$
。$\textit{Algebraic-Stacks}'/S$ 的对象 $\mathcal{X}'$
等价于某个 $\textit{Algebraic-Stacks}/S$ 中的 $\mathcal{X}$
所对应的 $f^{-1}\mathcal{X}$，当且仅当它有一个表现
% P11-E0077：冻结源此处写成 \mathcal{X}=[U'/R']，上下文对象为 \mathcal{X}'。
$\mathcal{X} = [U'/R']$，其中 $U', R'$ 分别同构于某些
$U, R \in \textit{Spaces}/S$ 所对应的 $f^{-1}U$、$f^{-1}R$。
\end{lemma}

\begin{proof}
关于态射范畴的陈述是更一般的《叠》引理
\ref{stacks-lemma-bigger-site} 的推论。“本质像”的刻画则由引理
\ref{algebraic-lemma-change-big-site} 的证明中对 $f^{-1}$ 的描述得出。
\end{proof}


\section{更换基概形}
\label{algebraic-section-change-base-scheme}

\noindent
本节简要讨论更换基概形时会发生什么。结论是：给定基概形之间的一个态射
$S \to S'$，任意 $S$ 上的代数叠都可以视为 $S'$ 上的代数叠。

\begin{lemma}
\label{algebraic-lemma-category-of-spaces-over-smaller-base-scheme}
设 $\Sch_{fppf}$ 是一个大 fppf 位点，并设
$S \to S'$ 是该位点中的一个态射。上文
Stacks, Section \ref{stacks-section-localize}
中的构造 A 与构造 B 给出如下 $2$-范畴的同构：
$$
\left\{
\begin{matrix}
S\text{ 上代数叠 }\mathcal{X}\\
\text{所成的 }2\text{-范畴}
\end{matrix}
\right\}
\leftrightarrow
\left\{
\begin{matrix}
\text{二元组 }(\mathcal{X}', f)\text{ 所成的 }2\text{-范畴，其中}\\
\mathcal{X}'\text{ 是 }S'\text{ 上的代数叠，而}\\
f : \mathcal{X}' \to (\Sch/S)_{fppf}\text{ 是 }S'\text{ 上代数叠的态射}
\end{matrix}
\right\}
$$
\end{lemma}

\begin{proof}
该陈述是有意义的，因为函子
$j : (\Sch/S)_{fppf} \to (\Sch/S')_{fppf}$
是与 $(\Sch/S')_{fppf}$ 的对象 $S/S'$ 相联系的局部化函子。由
Stacks, Lemma \ref{stacks-lemma-localize-stacks}
只需证明构造 A 与构造 B 保持代数叠的子范畴。例如，若
$\mathcal{X} = [U/R]$，则把构造 A 应用于 $\mathcal{X}$ 所得到的正是
$\mathcal{X}' = \mathcal{X}$。反之，若 $\mathcal{X}' = [U'/R']$，
则态射 $p$ 诱导代数空间的态射 $U' \to S$ 与 $R' \to S$；于是
$\mathcal{X} = [U'/R']$，但现在将其视为 $S$ 上的叠。因此引理成立。
\end{proof}

\begin{definition}
\label{algebraic-definition-viewed-as}
设 $\Sch_{fppf}$ 是一个大 fppf 位点，并设
$S \to S'$ 是该位点中的一个态射。若
$p : \mathcal{X} \to (\Sch/S)_{fppf}$
是 $S$ 上的代数叠，则所谓把 $\mathcal{X}$
{\it 视为 $S'$ 上的代数叠}，是指代数叠
$$
\mathcal{X} \longrightarrow (\Sch/S')_{fppf}
$$
它由把引理 \ref{algebraic-lemma-category-of-spaces-over-smaller-base-scheme}
中的构造 A 应用于 $\mathcal{X}$ 而得到。
\end{definition}

\noindent
反过来，若从 $S'$ 上的一个代数叠 $\mathcal{X}'$ 出发，并希望得到
$S$ 上的代数叠，则考虑 $2$-纤维积
$$
\mathcal{X}'_S
=
(\Sch/S)_{fppf} \times_{(\Sch/S')_{fppf}} \mathcal{X}'
$$
由引理 \ref{algebraic-lemma-2-fibre-product}，它是 $S'$ 上的代数叠。此外，
它还配有自然的 $1$-态射
$p : \mathcal{X}'_S \to (\Sch/S)_{fppf}$，因而由
引理 \ref{algebraic-lemma-category-of-spaces-over-smaller-base-scheme}，
它以典范方式对应于 $S$ 上的一个代数叠。

\begin{definition}
\label{algebraic-definition-change-of-base}
设 $\Sch_{fppf}$ 是一个大 fppf 位点，设
$S \to S'$ 是该位点中的一个态射，并设 $\mathcal{X}'$ 是
$S'$ 上的一个代数叠。所谓 $\mathcal{X}'$ 的{\it 基变换}，
是指上面所述的 $S$ 上代数叠 $\mathcal{X}'_S$。
\end{definition}
